%! TEX root = **/000-main.tex
% vim: spell spelllang=en:

%%%%%%%%%%%%%%%%%%%%%%%%%%%%%%%%%%%%%%%%%%%%%%%%%%%%%%%%%%%%%%%%%%%%%%%%%%%%%%%%
% PREAMBLE
%%%%%%%%%%%%%%%%%%%%%%%%%%%%%%%%%%%%%%%%%%%%%%%%%%%%%%%%%%%%%%%%%%%%%%%%%%%%%%%%
%! TEX root = **/000-main.tex
%%%%%%%%%%%%%%%%%%%%%%%%%%%%%%%%%%%%%%%%%%%%%%%%%%%%%%%%%%%%%%%%%%%%%%%%%%%%%%%%
%% LaTeX preamble, load in main.tex with: \input{preamble}
%%%%%%%%%%%%%%%%%%%%%%%%%%%%%%%%%%%%%%%%%%%%%%%%%%%%%%%%%%%%%%%%%%%%%%%%%%%%%%%%

\documentclass[12pt, oneside]{book}
\usepackage[a4paper, left=2.5cm, right=2.5cm, top=2.5cm, bottom=2.5cm]{geometry}

% for debugging overfulls
%\documentclass[draft, 12pt, oneside]{article}
%\usepackage[showframe, a4paper, left=2.5cm, right=2.5cm, top=2.5cm, bottom=2.5cm]{geometry}

%%%%%%%%%%%%%%%%%%%%%%%%%%%%%%%%%%%%%%%%%%%%%%%%%%%%%%%%%%%%%%%%%%%%%%%%%%%%%%%%
%% FONTS
%%%%%%%%%%%%%%%%%%%%%%%%%%%%%%%%%%%%%%%%%%%%%%%%%%%%%%%%%%%%%%%%%%%%%%%%%%%%%%%%

\usepackage[T1]{fontenc}
\usepackage{fontspec}
\usepackage{microtype}

\setmonofont[Scale=MatchLowercase]{DejaVu Sans Mono}

%%%%%%%%%%%%%%%%%%%%%%%%%%%%%%%%%%%%%%%%%%%%%%%%%%%%%%%%%%%%%%%%%%%%%%%%%%%%%%%%
%% LANGUAGE
%%%%%%%%%%%%%%%%%%%%%%%%%%%%%%%%%%%%%%%%%%%%%%%%%%%%%%%%%%%%%%%%%%%%%%%%%%%%%%%%

\usepackage{polyglossia}
\setdefaultlanguage{english}
\setotherlanguages{spanish,catalan}

%%%%%%%%%%%%%%%%%%%%%%%%%%%%%%%%%%%%%%%%%%%%%%%%%%%%%%%%%%%%%%%%%%%%%%%%%%%%%%%%
%% BIBLIOGRAPHY
%%%%%%%%%%%%%%%%%%%%%%%%%%%%%%%%%%%%%%%%%%%%%%%%%%%%%%%%%%%%%%%%%%%%%%%%%%%%%%%%

\usepackage[
    backend=biber,
    style=numeric,
]{biblatex}
\DeclareNameAlias{default}{family-given}

\addbibresource{biblio.bib}

\usepackage{fvextra}        % Req by minted (must load before csquotes)
\usepackage{csquotes}       % For bibliography quotations
\DeclareQuoteAlias{spanish}{catalan}

%%%%%%%%%%%%%%%%%%%%%%%%%%%%%%%%%%%%%%%%%%%%%%%%%%%%%%%%%%%%%%%%%%%%%%%%%%%%%%%%
%% COMMON
%%%%%%%%%%%%%%%%%%%%%%%%%%%%%%%%%%%%%%%%%%%%%%%%%%%%%%%%%%%%%%%%%%%%%%%%%%%%%%%%

\usepackage{color}     % more colors
\usepackage[dvipsnames]{xcolor}


\usepackage{graphicx}   % graphics
\graphicspath{{./figures/}}

\usepackage{comment}

%%%%%%%%%%%%%%%%%%%%%%%%%%%%%%%%%%%%%%%%%%%%%%%%%%%%%%%%%%%%%%%%%%%%%%%%%%%%%%%%
%% MATHS
%%%%%%%%%%%%%%%%%%%%%%%%%%%%%%%%%%%%%%%%%%%%%%%%%%%%%%%%%%%%%%%%%%%%%%%%%%%%%%%%

\usepackage{mathtools}  % amsmath + more
\usepackage{amsthm}     % Theorem enviroment
\usepackage{amssymb}    % More symbols
\usepackage{amstext}    % Text inside mathenv

%\usepackage{relsize}    % Bigger math with mathlarger{___}
%\usepackage{nicefrac}   % nice fractions in one line

%\usepackage{IEEEtrantools} % Complex equation arrays

%%%%%%%%%%%%%%%%%%%%%%%%%%%%%%%%%%%%%%%%%%%%%%%%%%%%%%%%%%%%%%%%%%%%%%%%%%%%%%%%
%% REFERENCES (load order is important)
%%%%%%%%%%%%%%%%%%%%%%%%%%%%%%%%%%%%%%%%%%%%%%%%%%%%%%%%%%%%%%%%%%%%%%%%%%%%%%%%

\usepackage{imakeidx}
\usepackage{varioref} % reference far away (1)
\usepackage[colorlinks = true]{hyperref} % links in references (2)
\usepackage{cleveref} % smart references (3)
%hyperref configuration so that it doesn't contrast so much colorlinks,
\hypersetup{
   linkcolor={black},
   citecolor={black},
   %linkcolor={red!50!black},
   %citecolor={blue!50!black},
   urlcolor={blue!80!black}
}

\usepackage[bottom]{footmisc} % Footnotes at bottom of page

%%%%%%%%%%%%%%%%%%%%%%%%%%%%%%%%%%%%%%%%%%%%%%%%%%%%%%%%%%%%%%%%%%%%%%%%%%%%%%%%
%% FIGURES
%%%%%%%%%%%%%%%%%%%%%%%%%%%%%%%%%%%%%%%%%%%%%%%%%%%%%%%%%%%%%%%%%%%%%%%%%%%%%%%%

%\usepackage[export]{adjustbox}  % Adjust table size
\usepackage{float}               % Force tables and images position (H and H!)
%\usepackage{wrapfig}            % Wrap images like in HTML

\usepackage[justification=centering]{caption}
%\usepackage{subcaption}                     % Subfigures
%\usepackage[framemethod=tikz]{mdframed}     % Custom frames

%%%%%%%%%%%%%%%%%%%%%%%%%%%%%%%%%%%%%%%%%%%%%%%%%%%%%%%%%%%%%%%%%%%%%%%%%%%%%%%%
%% TABLES
%%%%%%%%%%%%%%%%%%%%%%%%%%%%%%%%%%%%%%%%%%%%%%%%%%%%%%%%%%%%%%%%%%%%%%%%%%%%%%%%

%\usepackage{colortbl, booktabs} % Better tables
%\usepackage{tabularx}
%\usepackage{longtable} % Multiple page table (does not work with tabularx)
\usepackage{xltabular, colortbl, booktabs} % longtable + tabularx (has bug with booktabs: fix below)

% Split cell in lines and more formating options inside table
\usepackage{array, multirow, multicol, makecell}

%%%
% bug fix for booktabs + xltabular incompatibility
\makeatletter
\def\@BTrule[#1]{%
  \ifx\longtable\undefined
    \let\@BTswitch\@BTnormal
  \else\ifx\hline\LT@hline
    \nobreak
    \let\@BTswitch\@BLTrule
  \else
     \let\@BTswitch\@BTnormal
  \fi\fi
  \global\@thisrulewidth=#1\relax
  \ifnum\@thisruleclass=\tw@\vskip\@aboverulesep\else
  \ifnum\@lastruleclass=\z@\vskip\@aboverulesep\else
  \ifnum\@lastruleclass=\@ne\vskip\doublerulesep\fi\fi\fi
  \@BTswitch}
\makeatother
%%%

%%%%%%%%%%%%%%%%%%%%%%%%%%%%%%%%%%%%%%%%%%%%%%%%%%%%%%%%%%%%%%%%%%%%%%%%%%%%%%%%
%% SIUNITX
%%%%%%%%%%%%%%%%%%%%%%%%%%%%%%%%%%%%%%%%%%%%%%%%%%%%%%%%%%%%%%%%%%%%%%%%%%%%%%%%

%\usepackage[alsoload=hep]{siunitx}          % SI units and uncertainties
%\sisetup{locale = FR}                       % Commas and so on for spanish
%\sisetup{separate-uncertainty=true}
%\sisetup{
  %per-mode=fraction,
  %fraction-function=\nicefrac
%}

%%%%%%%%%%%%%%%%%%%%%%%%%%%%%%%%%%%%%%%%%%%%%%%%%%%%%%%%%%%%%%%%%%%%%%%%%%%%%%%%
%% TIKZ
%%%%%%%%%%%%%%%%%%%%%%%%%%%%%%%%%%%%%%%%%%%%%%%%%%%%%%%%%%%%%%%%%%%%%%%%%%%%%%%%

%\usepackage{tikz}
%\usetikzlibrary{arrows}
%\usetikzlibrary{scopes}
%\usetikzlibrary{babel}

%%%%%%%%%%%%%%%%%%%%%%%%%%%%%%%%%%%%%%%%%%%%%%%%%%%%%%%%%%%%%%%%%%%%%%%%%%%%%%%%
%% MINTED
%%%%%%%%%%%%%%%%%%%%%%%%%%%%%%%%%%%%%%%%%%%%%%%%%%%%%%%%%%%%%%%%%%%%%%%%%%%%%%%%

\usepackage{minted}
\definecolor{codeBg}{HTML}{FFFDE7}
\setminted{
    %style=pastie,
    frame=lines,
    framesep=3mm,
    linenos,
    breaklines=true,
    encoding=utf8,
    fontsize=\footnotesize,
    bgcolor=codeBg
}

%%%%%%%%%%%%%%%%%%%%%%%%%%%%%%%%%%%%%%%%%%%%%%%%%%%%%%%%%%%%%%%%%%%%%%%%%%%%%%%%
%% CUSTOM COMMANDS
%%%%%%%%%%%%%%%%%%%%%%%%%%%%%%%%%%%%%%%%%%%%%%%%%%%%%%%%%%%%%%%%%%%%%%%%%%%%%%%%

% empty whitepage without numbering
\newcommand{\whitepage}{
    \clearpage\thispagestyle{empty}\addtocounter{page}{-1} \newpage \clearpage
}

% Add command before appendix section for page numbering: A-1
\newcommand{\appendixpagenumbering}{
    \break
    \pagenumbering{arabic}
    \renewcommand{\thepage}{\thesection-\arabic{page}}
}

%%%%%%%%%%%%%%%%%%%%%%%%%%%%%%%%%%%%%%%%%%%%%%%%%%%%%%%%%%%%%%%%%%%%%%%%%%%%%%%%
%% CUSTOM MATH OPERATORS (functions not in italic in math mode)
%%%%%%%%%%%%%%%%%%%%%%%%%%%%%%%%%%%%%%%%%%%%%%%%%%%%%%%%%%%%%%%%%%%%%%%%%%%%%%%%

%\DeclareMathOperator{\arcsec}{arcsec}
%\DeclareMathOperator{\arccot}{arccot}
%\DeclareMathOperator{\arccsc}{arccsc}
%\DeclareMathOperator{\cis}{cis}

%%%%%%%%%%%%%%%%%%%%%%%%%%%%%%%%%%%%%%%%%%%%%%%%%%%%%%%%%%%%%%%%%%%%%%%%%%%%%%%%
%% MISC
%%%%%%%%%%%%%%%%%%%%%%%%%%%%%%%%%%%%%%%%%%%%%%%%%%%%%%%%%%%%%%%%%%%%%%%%%%%%%%%%

%\usepackage{datetime} % Customize date
%% \monthyeardate\today gives the date without the day
%\newdateformat{monthyeardate}{%
    %\monthname[\THEMONTH], \THEYEAR}

% Set default figure size
\setkeys{Gin}{width=.9\textwidth, keepaspectratio}

\usepackage{fancyhdr}
\usepackage{relsize}

% Center figures by default
\makeatletter
\g@addto@macro\@floatboxreset\centering
\makeatother

\usepackage[normalem]{ulem}
\usepackage{dsfont}
\usepackage{tikz}
% \usepackage{tikzexternal}
\usepackage{pgfplots}
\pgfplotsset{compat=1.18}

\usetikzlibrary{shapes,arrows,positioning,calc}
\usetikzlibrary{arrows.meta}
\tikzset{>={Latex[width=3mm,length=3mm]}}

\usepackage{cancel}
\usepackage{mathrsfs}

\DeclareMathOperator{\argmin}{argmin}
\DeclareMathOperator{\argmax}{argmax}

\directlua{require("lua/hyphenate_all")}


%%%%%%%%%%%%%%%%%%%%%%%%%%%%%%%%%%%%%%%%%%%%%%%%%%%%%%%%%%%%%%%%%%%%%%%%%%%%%%%%
% EXTRA PACKAGES / CONFIG
%%%%%%%%%%%%%%%%%%%%%%%%%%%%%%%%%%%%%%%%%%%%%%%%%%%%%%%%%%%%%%%%%%%%%%%%%%%%%%%%

% Set default figure size
\setkeys{Gin}{width=.9\textwidth, keepaspectratio}

\usepackage{fancyhdr}
\usepackage{relsize}

% Center figures by default
\makeatletter
\g@addto@macro\@floatboxreset\centering
\makeatother

\usepackage[normalem]{ulem}
\usepackage{dsfont}
\usepackage{tikz}
% \usepackage{tikzexternal}
\usepackage{pgfplots}
\pgfplotsset{compat=1.18}

\usetikzlibrary{shapes,arrows,positioning,calc}

\usepackage{cancel}
\usepackage{mathrsfs}

\DeclareMathOperator{\argmin}{argmin}
\DeclareMathOperator{\argmax}{argmax}

\usepackage{tcolorbox}
\tcbuselibrary{most}

\theoremstyle{plain}% from `amsthm'
% \newtheorem{lemma}{Lemma}% from `amsthm'
% \newtheorem{prop}{Proposition}% from `amsthm'

\theoremstyle{definition}% from `amsthm'
\newtheorem{defn}{Definition}% from `amsthm'

\theoremstyle{remark}% from `amsthm'
\newtheorem*{rem}{Remark}% from `amsthm'
\newtheorem{ex}{Example}% from `amsthm'
\newtheorem{prob}{Problem}% from `amsthm'
\newtheorem{ques}{Question}% from `amsthm'
\newtheorem*{note}{Note}% from `amsthm'
\newtheorem*{hint}{Hint}% from `amsthm'


\tcbset{%
    theo/.style={%
        enhanced,
        skin=bicolor,
        colbacklower=#1!2,
        sharp corners,
        toprule=0pt, rightrule=0pt, bottomrule=0pt, leftrule=1mm,
        colback=#1!5, colframe=#1!80!black, coltitle=#1!80!black,
        fonttitle=\bfseries,
        detach title, before upper={\tcbtitle\quad},
        subtitle style={fontupper=\bfseries,toprule=0pt,bottomrule=0pt,colback={#1!10}},
    },
    marker/.style={%
      enhanced,
      before skip balanced=2mm,after skip balanced=3mm,
      boxrule=0.4pt,left=5mm,right=2mm,top=1mm,bottom=1mm,
      colback=#1!50,
      colframe=#1!20!black,
      sharp corners,rounded corners=southeast,arc is angular,arc=3mm,
      underlay={%
        \path[fill=tcbcolback!80!black] ([yshift=3mm]interior.south east)--++(-0.4,-0.1)--++(0.1,-0.2);
        \path[draw=tcbcolframe,shorten <=-0.05mm,shorten >=-0.05mm] ([yshift=3mm]interior.south east)--++(-0.4,-0.1)--++(0.1,-0.2);
        \path[fill=#1!50!black,draw=none] (interior.south west) rectangle node[white]{\Huge\bfseries !} ([xshift=4mm]interior.north west);
        },
      % drop fuzzy shadow,
  },
}

\newtcbtheorem[number within=chapter]{definition}{Definition}%
{theo=blue}{def}

\newtcbtheorem[number within=chapter]{theorem}{Theorem}%
{theo=purple}{th}

\newtcbtheorem[number within=chapter]{prop}{Proposition}%
{theo=red}{prop}

\newtcbtheorem[number within=chapter]{lemma}{Lemma}%
{theo=magenta}{lemma}

\newtcbtheorem[number within=chapter]{recap}{Recap}%
{theo=orange}{rec}

\newtcbtheorem[number within=chapter]{exercise}{Exercise}%
{theo=yellow,coltitle=black}{exer}

\newtcbtheorem[number within=chapter]{question}{Question}%
{theo={green!50!black}}{qu}

\newtcbtheorem[number within=chapter]{example}{Example}%
{theo=black,colback=white}{ex}

\newtcolorbox{marker}[1][]{marker=yellow,#1}
\newtcolorbox{important}[1][]{marker=red,#1}
\newtcolorbox{cnote}[1][]{marker=orange,#1}

%%%%%%%%%%%%%%%%%%%%%%%%%%%%%%%%%%%%%%%%%%%%%%%%%%%%%%%%%%%%%%%%%%%%%%%%%%%%%%%%
% METADATA
%%%%%%%%%%%%%%%%%%%%%%%%%%%%%%%%%%%%%%%%%%%%%%%%%%%%%%%%%%%%%%%%%%%%%%%%%%%%%%%%

% remove when using \maketitle:
\renewcommand\and{\\[\baselineskip]}

\title{\Huge Advanced Machine Learning\\ \Large Lecture Notes}
\author{Aleix Boné}
\date{\today}

\makeindex

\begin{document}
\newcommand{\iemph}[1]{\index{#1}\emph{#1}}
%%%%%%%%%%%%%%%%%%%%%%%%%%%%%%%%%%%%%%%%%%%%%%%%%%%%%%%%%%%%%%%%%%%%%%%%%%%%%%%%
% TITLE
%%%%%%%%%%%%%%%%%%%%%%%%%%%%%%%%%%%%%%%%%%%%%%%%%%%%%%%%%%%%%%%%%%%%%%%%%%%%%%%%

% Default title or use titlepage.tex

%\maketitle
%! TEX root = **/000-main.tex
% vim: spell spelllang=en:

\thispagestyle{empty}
\clearpage
\setcounter{page}{-1}

\makeatletter
\begin{titlepage}
{
    \centering
    % \includegraphics[width=0.9\textwidth]{institution-logo}
    \null%
    % \vspace{3em}
    \vspace{13em}
    {\Huge \bfseries \@title \\[1em]
      \Large \scshape Lecture Notes
    \par}
    \vspace{2em}
    {\large \scshape \@date \par}
    \vfill
\begin{center}
    % Supplementary image
\end{center}
    \vspace{5em}

    \vfill
    {\raggedleft \large Notes Author: \@author \par}
    % {\raggedleft \large Instructor: Lluís Belanche \par}
}
\end{titlepage}
\makeatother


%%%%%%%%%%%%%%%%%%%%%%%%%%%%%%%%%%%%%%%%%%%%%%%%%%%%%%%%%%%%%%%%%%%%%%%%%%%%%%%%
% TOC & lists
%%%%%%%%%%%%%%%%%%%%%%%%%%%%%%%%%%%%%%%%%%%%%%%%%%%%%%%%%%%%%%%%%%%%%%%%%%%%%%%%

\pagenumbering{Roman}

%\setcounter{tocdepth}{2}
\tableofcontents \pagebreak

\listoffigures \listoftables \clearpage

\pagenumbering{arabic}

%%%%%%%%%%%%%%%%%%%%%%%%%%%%%%%%%%%%%%%%%%%%%%%%%%%%%%%%%%%%%%%%%%%%%%%%%%%%%%%%
% SECTIONS
%%%%%%%%%%%%%%%%%%%%%%%%%%%%%%%%%%%%%%%%%%%%%%%%%%%%%%%%%%%%%%%%%%%%%%%%%%%%%%%%

% Paragraph spacing (placed after ToC)
\setlength{\parskip}{1em plus 0.5em minus 0.2em}
%\setlength{\parindent}{0pt}

\setlength{\headheight}{14.5pt}
\pagestyle{fancy}

% Theory
\part{Theory}
\section{Setting of the learning problem}

\subsection{Learning problem definition}

Can be defined as 3 components:
\begin{itemize}
	\item
	      A \emph{generator} \(G\) of random data \(x\) from a
	      probability distribution \(P(x)\)which is fixed and unknown to us.
	\item
	      A \emph{supervisor} \(S\) which returns an output value ${y \in Y}$
	      for each everyinput \(x \in X\) according to a conditional
	      probability distribution \(P(y | x)\), also fixed and unknown
	      to us.
	\item
	      A \emph{learning machine (algorithm)} \(L\) which is capable of
	      implementing(computing) any of a set of functions (models).
	      \(f_\theta\) where \(\theta\) is a set of parameters belonging to the
	      parameter space \(\theta\).
\end{itemize}

\begin{center}
	\tikzstyle{block} = [draw, rectangle, minimum height=3em, minimum width=3em]
	\tikzstyle{virtual} = [coordinate]
	\begin{tikzpicture}[>=stealth, auto, node distance = 1cm and 3cm]
		% Place nodes
		\node [block] (G) {$G$};
		\node [block, right=of G] (S) {$S$};
		\node [block, below=of S] (L) {$L$};
		\node [virtual, right=of S] (So) {};
		\node [virtual, right=of L.340] (Lo) {};
		\node [below right=of L, align=center] (d) {$L(y,\,\ddot{y})$ \\ (discrepancy)};

		\draw [->] (G.east) -- node [name=x] {$x$} (S);
		\draw [->, rounded corners] (x) |- (L);

		\draw [->] (S) -- node [name=y] {$y$} (So);
		\draw [->, rounded corners] (y) |- (L.20);

		\draw [->] (L.340) -- node[below] {$\ddot{y}$} (Lo);

		\draw [->, rounded corners] (d.west) -| node[below] {modify} (L.south);
	\end{tikzpicture}
\end{center}

The selection of the best possible model according to \(L\) is done
based on a learning data sample
\(D = \bigl\{(x_1,\, y_1),\, (x_2,\, y_2),\, ...,\, (x_n,\, y_n)\bigr\}\) (where \(x_i\) are
vectors)

\subsection{Why is the relation between X and Y stochastic?}

In most of the cases, randomness is ignorance of the true effects.
\begin{enumerate}
	\item
	      Ignorance of the functional dependence on a not measured variable
	      \(z\).
	\item
	      Technical invisibility.
\end{enumerate}

\subsection{Some definitions}

\begin{definition}{Loss function}{}
	A function
	\(L: Y \times Y \longrightarrow \mathds{R}+\) such that:

	\begin{align*}
		L(y,\, y') & = L(y',\, y) \tag{symmetry} \\
		L(y,\, y') & = 0  \Longleftarrow y = y'
	\end{align*}

	Examples of loss functions:

	\begin{align*}
		L(y,\, y') & := \begin{cases}1 \quad \text{if } y \neq y' \\ 0 \quad \text{otherwise} \end{cases} &  & \text{(0-1 loss) / (zero-one) / } L_{01} \\
		L(y,\, y') & := (y - y')^2                                                                        &  & \text{(squared error)}
	\end{align*}
\end{definition}

\begin{definition}{Hypothesis space}{}
	\begin{equation*}
		\mathcal{F} = \bigl\{f_\theta: \mathcal{X} \rightarrow \mathcal{Y} \mid \theta \in \varTheta \subseteq \mathds{R} \bigr\}
	\end{equation*}
\end{definition}


\begin{definition}{Classifier}{}
	Is a function \(f: X \longrightarrow Y\)
	where \(Y\) is a set of different symbols.
\end{definition}

\begin{definition}{Risk}{}

	The risk of a function is its \emph{expected loss}. Also known as true
	error or generalization error.

	\begin{equation*}
		\mathds{R}[f] := \mathds{E}_{(x,\, y) \sim p} \left[ L(f(x),\, y) \right]
	\end{equation*}

	\begin{note}
		A function that takes a function as an argument is called a functional.
		Risk is a functional.
	\end{note}
\end{definition}

\begin{definition}{Empirical error}{}

	The average error or loss function evaluated on a specific data sample
	\(D\).

	\begin{equation*}
		\hat{R}_{D^n}[f] = \frac{1}{n} \sum_{i=1}^n L(f(x_i), y_i)
	\end{equation*}

	\begin{note}
		This is also a functional.
	\end{note}
\end{definition}

\begin{definition}{Training error}{}
	The \emph{empirical error} obtained in the data
	sample used for training.
\end{definition}

\begin{prop}{Model error proposition}{}

	For a fixed model \(f_θ\) the expected value of the
	empirical error based on a data sample is equal to the true error.

	\begin{equation*}
		\mathds{E}_{D^{ijd}p^n}[\hat{R}_{D^n}[f]] = R[f]
	\end{equation*}

	\begin{proof} TODO
	\end{proof}
\end{prop}


\subsection{Discriminative vs generative classifiers}

\begin{definition}{Discriminative classifier}{discriminative}
	A classifier that directly models the conditional probability
	\(P(y | x)\) of the output given the input.

	\paragraph{Example:} Logistic regression
\end{definition}

\begin{definition}{Generative classifier}{generative}
	A classifier that models the joint probability distribution
	\(P(x, y)\) of the input and the output.

	\paragraph{Example:} Naïve Bayes
\end{definition}

\subsection{Risk minimizers}

\begin{align*}
	R[f] := \mathds{E}_{(x, y) \sim p} [ L(f(x), y) ] \\
	\int_{\mathds{R}^d} \int_{\mathds{R}} L(f(x), y) p(x, y)\; dy\, dx
\end{align*}

\begin{prop}{}{}

	The function minimizing the risk on the zero-one loss is the Bayes
	classifier. Which is the one that given x, assigns the class with the
	highest posterior probability.

	\[f^* = \argmax {P(w | x)}\]

	\begin{enumerate}
		\item The risk of \(f_{01}^*\) is the Bayes risk.
		\item The upper bound of the Bayes risk is 0.5
	\end{enumerate}

\end{prop}

\subsection{Regression under the square error / loss function}

\subsubsection{Standard setting of the regression problem}

\begin{align*}
	y & = f_\text{sq.}(x,\, z)                         \\
	  & \downarrow (z \text{ is absent from the data}) \\
	y & = f_\text{sq.}(x) + \varepsilon
\end{align*}

Where:
\begin{itemize}
	\item \(\mathds{E}[\varepsilon] = 0\) (expected value of the error is 0)
	\item \(\text{Var}[\varepsilon] = \sigma^2 < \infty\) (variance of the error is finite) (a.k.a. homoscedasticity)
	\item No assumption of gaussianity
\end{itemize}

\begin{prop}{}{}
	\begin{align*}
		R[f]                       & = \sigma^2 + \int_{\mathds{R}^d} (f(x) - f^*_{sq.}(x))^2 p(x)\; dx
		                           &                                                                    & (\text{MSE}[f])              \\
		\text{where } f^*_{sq.}(x) & = \int_{\mathds{R}} y\cdot p(y(x))\; dy
		                           &                                                                    & \text{(regression function)} \\
		\text{in terms of statistics: }                                                                                                \\ \mathds{E}_{Y|x}[y | X = x] \\
		\Rightarrow R[f^*_{sq.}]
	\end{align*}

	We cannot compute this in practice since we do not know the distribution
	of the data.
\end{prop}

\begin{prop}{bias-variance decomposition}{}

	The bias / variance decomposition of the mean squared error.

	\[MSE[f] = Bias(f)^2 + Var(f)
	\]

	There is a trade-off between bias and variance, known as the
	\emph{bias-variance dilemma}.
\end{prop}

\paragraph{Remainder}

Given a dataset \(D = {x^i, y^i}\) of size \(n\) \((x^i, y^i)\sim p\) We
choose a loss function \(L\) and a hypothesis space \(F\).

\[F := { x \rightarrow f_\theta(x),\quad \theta \in \Theta }
\]

\subparagraph{Observations}

The set of functions we chose matters a lot. In some sense \(F\) should
be large to minimize the chances that the chosen function is not in
\(F\).

\begin{itemize}
	\item
	      Choosing a Machine Learning Algorithm of \(F\) given a particular
	      \(D_n\) obviously depends on this \(D_n\):

	      \[D_n! \rightsquigarrow f_\theta = f_n\]
	\item The best possible solution \(f^*_{sq.}\) may not belong to \(F\).

	      In case \(f^*_{sq.} \notin F\) then we should find:
	      \(\hat{f}_{sq.} = \argmin {\lVert f - f^*_{sq.} \rVert}\)
	      We could use different norm functions.
\end{itemize}

The expression being integrated is:

\[\left( f_m(x) - f^*_{sq.}(x) \right)^2
\]

Expected value of the model for all possible datasets in \(D_n\):
\begin{align*}
	\mathds{E}_{D_n} \left[ f_n(x) \right] & =: \overline{f_n}                                               \\
	\left( \left( f_n(x) - \overline{f_n} \right) + \left( \overline{f_n} - f_{sq.}^*(x) \right) \right)^2 = \\
	\left( \overline{f_n}(x) - f_{sq.}^*(x) \right)^2 +
	\mathds{E}_{D_n} \left[ \left( f_n(x) - \overline{f_n}(x) \right)^2 \right]
	\rightarrow \text{Var}(f_n(x))
\end{align*}
How the prediction coming from a specific data sample $D_n$ varies around its mean
(Dependence of the model being learned on the specific data sample)

\(\text{Bias}^2(f_n(x))\) : how the average prediction over all possible
\(D_n\) at point \(x\) differs from the best possible prediction.

\begin{align*}
	\text{Bias}^2(f) & = \int_{\mathds{R}^d} \left( \overline{f_n}(x) - f_{sq.}^*(x) \right)^2 p(x)\; dx                           \\
	\text{Var}(f)    & = \int_{\mathds{R}^d} \mathds{E}_{D_n} \left[ \left( f_n(x) - \overline{f_n}(x) \right)^2 \right] p(x)\; dx \\
	\text{MSE}[f]    & = Bias(f)^2 + Var(f) + \sigma^2
\end{align*}

where \(\sigma^2\) is the variance of the noise.

Since \(\sigma\) cannot be reduced, it is called irreducible error. The
rest is called reducible error.

Getting better data we reduce the irreducible error. (ignorance)

\section{Consistency}

We depart from a specific data sample \(D_n\) The training error of a
model \(f\) is:
\begin{equation*}
	\hat{R}_n^{[f]} := \frac{1}{n} \sum_{i=1}^n L(f(x^i),\, y^i)
\end{equation*}

So, when we learn a specific model \(f_n\) from \(D_n\) its
\emph{training error} is \(\hat{R}_n(f_n)\) The true error of \(f_n\)
is \(R(f_n)\)

The \emph{Empirical Risk Minimization} (ERM) prescribes to minimize the
training error.

\begin{description}
	\item[Consistency]
		The ERM is said to be \emph{consistent} if:

		\begin{align*}
			\hat{R}_n(f_n) & \xrightarrow[\text{infimum}_{f \in F}]{} R(f)
			\text{ as } n \rightarrow \infty                               \\
			R(f_n)         & \xrightarrow[\text{infimum}_{f \in F}]{} R(f)
			\text{ as } n \rightarrow \infty
		\end{align*}

		These convergences are in probability.
\end{description}

Graphically:
\begin{figure}[H]
	\begin{tikzpicture}
		\begin{axis}[
				domain=0.1:1.1,
				xmin=0.2, xmax=1,
				ymin=-6, ymax=6,
				samples=100,
				axis y line=center,
				axis x line=bottom,
				xlabel = {$n$},
				legend pos=south east,
				% area style,
				ticks=none,
			]
			\addlegendentry{\(R(f_n)\)};
			\addlegendentry{\(\hat{R}_n(f_n)\)};
			\addlegendentry{Convergence};
			\addplot+[mark=none] {1/x};
			\addplot+[mark=none] {-1/x};
			\addplot+[mark=none,dashed] {0};
			\node at (axis cs:1,1.1) [anchor=north east] {\(\text{infimum}_{f \in \mathcal{F}} R(f)\)};
		\end{axis}

	\end{tikzpicture}
	\caption{Convergence of the training error and the true error}
\end{figure}

\begin{theorem}{Vapnik-Chervonenkis}{}

	A necessary and sufficient condition for consistency is that the
	hypothesis space \(F\) is compact.

  % TODO: formal definition
	% \begin{equation*}
	% 	\forall \varepsilon > \dots
	% \end{equation*}
\end{theorem}
 % Problemes
\section{A glimpse into Bayesian methods (for supervised ML)}
\index{bayesian}

We have a data sample \(D = \{ x^i, y^i\}\quad i = 1, \dots, n\)
where \(x \in \mathcal{X} = \mathds{R}^d, y^i \in \mathcal{Y}\)
We consider a hypothesis space
\(\mathcal{F} := \{ f_\theta (x), \theta \in \Theta \subset \mathds{R}^d \}\).

The Bayesian idea is not to consider the optimal solution
\(\hat{\theta}\) as a point solution, but as a set of possible solutions,
some of which are more likely (to be correct) than others.

The available data \(D\) serves to reduce the uncertainty of the
distribution.

The Bayesian machinery makes use of the Bayes formula.

\begin{align*}
	P(D | \theta) \tag{likelihood} \\
	P(\theta | D) \tag{posterior}  \\
	P(\theta) \tag{prior}          \\
	\int_\Theta P(D | \theta') P(\theta') d\theta' = P(D) \tag{EXPECTED \equiv EVIDENCE}
\end{align*}

\[
	P(\theta | D) = \frac{P(D | \theta) P(\theta)}
	{ \int_\Theta P(D | \theta') P(\theta') d\theta' }
\]

\subsection{Procedure}

\begin{quote}
	Modelling is choosing
\end{quote}

\begin{enumerate}
	\item
	      Collect the data \(D\)

	      \begin{enumerate}
		      \item
		            Decide the \emph{functional form of the models} \(f_\theta (x)\)
		      \item
		            Choose the \emph{prior} \(P(\theta)\) about the parameters
	      \end{enumerate}
	\item
	      Calculate the \emph{likelihood function} \(P(y | x, \theta)\) (The
	      probability of getting the data \(D\) for a specific value of
	      \(\theta\))
	\item
	      Calculate the \emph{posterior distribution} of \(\theta\)

	      \[
		      P(\theta | D) = \frac{P(\theta) \prod_{i=1}^n}{P(y^i|x^i, \theta)
			      {\int_\Theta d\theta' P(\theta')\prod_{i=1}^n} P(y^i|x^i, \theta')}
	      \]

	\item
	      Decide a method for making predictions (when new data \(x^0\) comes):

	      \begin{enumerate}
		      \item Eliminate the prior \(P(\theta)\) \textrightarrow{}
		            \(P(\theta | D) \propto \mathcal{L}(\sigma ; D) \cdot P(\theta)\)

		            \begin{enumerate}
			            \item
			                  Then, the only thing that matters is the likelihood
			            \item
			                  You then find a set of parameters \(\theta\) that maximizes the
			                  likelihood function
			                  \(\hat{\theta}_{ML} = \argmax_{\theta \in \Theta} \mathcal{L}(\sigma ; D)\)
		            \end{enumerate}
		      \item
		            Half Bayesian
		            \(P(\theta | D) \propto \mathcal{L}(\sigma ; D) \cdot P(\theta)\)

		            \begin{enumerate}
			            \item
			                  \(\hat{\theta}_{MAP} = \argmax_{\theta \in \Theta} \mathcal{L}(\sigma ; D) \cdot P(\theta)\)
			                  (Maximum a posteriori)
			            \item
			                  Since we have a multiplication to minimize, we use logarithms.
			            \item
			                  It is then:
			                  \(\hat{\theta}_{MAP} = \argmax_{\theta \in \Theta} \log \mathcal{L}(\sigma ; D) + \log P(\theta)\)
		            \end{enumerate}
		      \item
		            Calculate the \emph{average of the posteriors} \(P(\theta | D)\)
                instead of the maximum. (average \textrightarrow{} expected value)

		            \begin{enumerate}
			            \item \(\hat{\theta}_{avg} = \int_\Theta \theta P(\theta | D) \; d\theta\)
			            \item In practice, this integral is intractable, we have to use
			                  numerical integration. Which is quite delicate.
		            \end{enumerate}
		      \item Fully Bayesian: you have to work with the full posterior.
		            \begin{enumerate}
			            \item You don't return any model. You return a distribution over models.
			            \item Since we have a distribution, we can compute confidence integrals
			                  (bayesian confidence intervals) (not in the scope of the course)
			            \item Suppose that a new data point comes, \(x^0\):

			                  \[P(y^0 | x^0, D) = \int_\Theta P(y^0 | x^0, \theta) P(\theta | D)\; d\theta
			                  \]

			                  In the first 2 cases, we obtain a single unique model.In all
			                  cases, the data is constant.
		            \end{enumerate}
	      \end{enumerate}
\end{enumerate}

\subsubsection{Fully worked example}

We depart from a data sample and face a regression task.

So \(x^i \in \mathds{R}^d,\; y^i \in \mathds{R}\) and we want to perform
\emph{linear regression}.

\begin{equation*}
	\mathcal{F}_{\text{linr}} := \bigl\{
	x \mapsto f_\omega(x) = \omega^T x + \omega_0 ,\quad
	\omega \in \mathds{R}^d, \quad
	\omega_0 \in \mathds{R}
	\bigr\}
\end{equation*}

By convention \(\underbar{\omega} := (\omega,\, \omega_0)\)

\begin{itemize}
	\item What loss function do we use?
	\item What kind of regularizer (penalty on model complexity) should I
	      use?
\end{itemize}

\begin{align*}
  y &= f_\omega(x) + \varepsilon \\
  \mathds{E}[\varepsilon] &= 0 \\
  \text{Var}(\varepsilon) &= \sigma^2 < \infty \\
  \text{Independence between }& x \text{ and }\varepsilon
\end{align*}

\paragraph{Procedure}

\begin{enumerate}
	\item For an arbitrary input \(x \in \mathds{R}^d\), we want to predict
	      \(y \in \mathds{R}\)
	      \begin{align*}
		      \epsilon                      & \sim \mathcal{N}(0, \sigma^2)                                       \\
		      P(y | x,\, \underbar{\omega})   & = \mathcal{N}(y | f_{\underbar{\omega}}(x),\, \sigma^2)               \\
		      \beta                         & = \frac{1}{\sigma^2} \text{ (precision)}                            \\
		      P(\underbar{\omega})          & = \mathcal{N}(\underbar{\omega} ;\, 0 ,\, \Sigma)                       \\
		      \Sigma                        & := \text{ covariance matrix (Not viable in practice, we use other)} \\
		      P(\underbar{\omega} | \alpha) & = \mathcal{N}(\underbar{\omega} ;\, 0 ,\,\sigma_{\omega}^2 I))         \\
		      \alpha                        & := \frac{1}{\sigma_{\omega}^2}
	      \end{align*}

	      A parameter that controls other parameters is called a
	      \iemph{hyperparameter}. In this case, \(\sigma_\omega\) is an
	      hyperparameter.

	\item The \emph{likelihood function}

	      \[
		      \mathcal{L}(\underbar{\omega} ; \beta) = \prod_{i=1}^n P(y^i | x^i, \underbar{\omega}) =
		      \prod_{i=1}^n \mathcal{N}(y^i | f_{\underbar{\omega}}(x^i), \sigma^2)
	      \]

	\item The \iemph{posterior distribution}
\end{enumerate}

\[P(\underbar{\omega} | y, X, \alpha, \beta) = \frac{\mathcal{L}(\underbar{\omega}; \beta) P(\underbar{\omega} | \alpha)}
	{\int_{\mathds{R}} \mathcal{L}(\underbar{\omega}; \beta) P(\underbar{\omega} | \alpha) d\underbar{\omega}}
\]

\begin{recap}{}{}

\subparagraph{Model}

\[
	P(y |  X, \underbar{\omega}, \beta) = \mathcal{N}(y | f_{\underbar{\omega}}(x), \beta) =
	\sqrt{\frac{\beta}{2\pi}} \exp \left( -\frac{\beta}{2} (y - f_{\underbar{\omega}}(x))^2 \right)
\]

\subparagraph{Prior}

\[P(\underbar{\omega} | \alpha) = \mathcal{N}(\underbar{\omega} ; 0 , \sigma_{\omega}^2 I)) =
	\sqrt{\frac{\alpha}{2\pi}} \exp \left( -\frac{\alpha}{2} \underbar{\omega}^T \underbar{\omega} \right)
\]

\subparagraph{Result}

The posterior distribution turns out to be:

\[P(\underbar{\omega} | D) = \mathcal{N}(\underbar{\omega} ; \mu, \Sigma)
	\mu = \beta \Sigma X^T y
	\Sigma = \left( \alpha I + \beta X^T X \right)^{-1}
\]

\begin{enumerate}
	\item
	      Maximum likelihood: \(\hat{\underbar{\omega}} = (X^T X)^{-1} X^T y\)
	\item
	      Maximum posteriori: \(\hat{\underbar{\omega}} = ???\) (Standard ridge
	      regression)
	\item
	      Average posteriori: \(\hat{\underbar{\omega}} = ???\)
	\item
	      Full Bayesian: \(\hat{\underbar{\omega}} = ???\)
\end{enumerate}

\end{recap}
 % Teoria
\subsection{Homework} % (fold)
\label{sub:Homework}

Ridge regression ($\lambda$)
\begin{equation}
	\rightsquigarrow \lambda ^* := \frac{.}{.}
\end{equation}
% subsection subsection name (end)

\section{Regularization-based learning} % (fold)
\label{sec:Regularization-based-learning}

Many \emph{MLA}s can be expressed as optimization problem.

\begin{itemize}
	\item We depart from:
	      \begin{itemize}
		      \item hypothesis space $\mathcal{F}$
		      \item loss function $L \geq 0$
		      \item a regularizer $\Omega \geq 0$ which expresses how to discourage complex models.
	      \end{itemize}
	\item
\[
	\min_{f \in \mathcal{F}} \left(\frac{1}{n} \sum_{i=1}^n L\left(f(x^i),\, y^i\right) + \lambda\Omega(f) \right)
	, \quad \lambda > 0 \text{ is a regularization parameter.}
\]

\end{itemize}

\subsection{Typical choices for the regularizer} % (fold)

\begin{align*}
	\Omega(f) & = ||f||_q^q, \quad q \geq 1                 \\
	\text{e.g. } q = 2 \text{ is the } \ell_2 \text{ norm.} \\
	||f||_2^2 & = \sqrt{\int f^2(x) \, dx}
\end{align*}

\subsection*{Is this problem convex?} % (fold)

\begin{itemize}
	\item All norms are convex.
	\item If $L$ is convex on its first argument $\Rightarrow$ the problem is convex.
\end{itemize}

Suppose we chose $\mathcal{F}, L, \Omega$ Questions:

\begin{itemize}
	\item How can I justify these choices?
	\item How do we optimize $\lambda \xrightarrow[D]{} \hat\lambda$?
	\item What kind of knowledge do we obtain upon solving the problem?
	\item How do we solve the general optimization problem?
\end{itemize}

\subsection{Learning models} % (fold)

\begin{align*}
	x \in \mathcal{X} = \mathbb{R}^d & , \quad y \in \mathcal{Y} = \mathbb{R}                                            \\
	f_{\text{basic model}}(x)        & = \omega_0 + \sum_{l = 1}^d \omega_l x_l, \quad \theta = \binom{\omega}{\omega_0}
\end{align*}

\paragraph{Linear model}

A model $f_\theta$ is linear when it has a linear dependence on its parameters.

ex:
\begin{align*}
	d      & = 1                            \\
	f_c(x) & = f(x;c)                       \\
	       & = \sum_{l=0}^d c_l \cdot (x)^l
\end{align*}

\subparagraph{Hint} Define $z_l = x^l$ and $c_l = \omega_l$.

\subsection{Homework}

Express a model in $\mathds{R}$ as a truncated Fourier series and
that it is a \emph{linear model}

\hrulefill

We can introduce:

\begin{equation*}
	\phi (x) = \left(\begin{array}{c}
			\phi_1(x) \\
			\vdots    \\
			\phi_n(x)
		\end{array}\right), \quad x \in \mathds{R}^d
\end{equation*}

Generalize on basic model:

\begin{equation}
	\boxed{f_\omega(x) = \sum_{l=1}^n \omega_l \phi_l(x)}
\end{equation}

A further generalization:

\begin{align*}
	f_\omega(x) = g\left(\sum_{l=1}^n \omega_l \phi_l(x)\right) = g\left(\omega^T\phi(x)\right) \\
	g: \mathds{R} \to \mathds{R} \text{ is a monotonic function (increasing)}                   \\
	\text{Typically non-linear}
\end{align*}

In the context of neural networks $g$ is called an \emph{activation function}.

For linear regression we have $g(z) = z$ (identity function) or
no activation function at all.

\paragraph{Examples of activation functions}

\begin{equation*}
	g(z) = \begin{cases}
		z & \text{if } z \geq \frac{1}{2} \\
		0 & \text{otherwise}
	\end{cases}  = \mathds{1}_{\{z \geq 1/2\}}(z), \quad z \in (0, 1)
\end{equation*}

We can set $\frac{1}{2}$ as a hyperparameter.

\subsection{A gentle introduction to GLMs}

Generalized Linear Models (GLMs) are a class of models that are
parametrized by a linear combination of features.

Recall:

\begin{align*}
	glm  & := g(\omega^T\phi(x))  \tag{\text{General Linear Model}} \\
	GMLs & \subset glms
\end{align*}

\newcommand\hcancel[2][black]{\setbox0=\hbox{$#2$}%
	\rlap{\raisebox{.45\ht0}{\textcolor{#1}{\rule{\wd0}{1pt}}}}#2}


\begin{align*}
	y^i         & = f_\omega(x^i) + \varepsilon^i, \quad \varepsilon^i \sim \mathcal{N}(0, \sigma^2) \\
	f_\omega(x) & = g\left(\omega^T\phi(x)\right) \quad \leftarrow \text{ glm}                       \\
	            & \hcancel[red]{\omega^T\phi(x) + \omega_0}                                          \\
	            & \Rightarrow y^i \sim \mathcal{N}\left(y; f_\omega(x^i), \sigma^2\right)            \\
	\mu_i       & := \mathds{E}[y^i  | X = x^i] = f_\omega(x^i) = \omega^T\phi(x^i)
\end{align*}

\begin{enumerate}
	\item Can we arrange for distributions other than the Gaussian?

	      It turns out that other distributions can be used. Those that
	      we name the \emph{exponential family of distributions}.
	      % \footnote{\href{https://en.wikipedia.org/wiki/Exponential_family#Table_of_distributions}{Wikipedia table}}

	      \begin{equation*}
		      P(y | \theta) = h(y)c(\theta) \exp\left(\sum_{j=1}^J \omega_j(\theta)t_j(x)\right)
	      \end{equation*}

	\item Can we merge the relation between our model's prediction
	      and the expected value of the predicted variable non-linear?

	      We admit a (normally) non-linear relation between $\mu_i$ and $f_\omega(x^i)$.

	      In the GLM context:

	      \begin{equation*}
		      \boxed{
			      h\left(\mathds{E} [y | X = x] \right) = \omega^T x
		      }
	      \end{equation*}

	      If $h$ admits an inverse function, then $g = h^{-1}$

	      The function $h$ is called the \emph{link function}.


\end{enumerate}

\paragraph{ex}.

\begin{center}
	\begin{tabular}{ccc}
		\toprule
		\text{Distribution}       & $g(z)$     & \text{Expected value} \\
		\midrule
		\text{Gaussian}           & $z$        & $\mu_i$               \\
		\text{Poisson}$(\lambda)$ & $\ln(z)$   & $\lambda$             \\
		\text{Bernoulli}$(p)$     & logit$(z)$ & $p$                   \\
		\bottomrule
	\end{tabular}
\end{center}

The last layer of a neural network is a GLM.

\subsection{Kinds of response/target variables}

\begin{center}
	\begin{tabular}{lll}
		\toprule
		variable     & regression  & distribution \\
		\midrule
		continuous   & linear      & gaussian     \\
		binary       & logistic    & bernoulli    \\
		nominal      & multinomial & multinomial  \\
		ordinal      & ordinal     & ordinal      \\
		count        & Poisson     & Poisson      \\
		failure time & survival    & Weibull      \\
		\bottomrule
	\end{tabular}
\end{center}

 % Teoria
\subsection{Procedure to learn a GLM}

Reference: \fullcite{dobson_introduction_2018}

\newcommand{\dcheckmark}{\checkmark~\hspace{-0.9em}\checkmark}

\begin{enumerate}
  \item Identify the kind of variable you want to predict. \checkmark
  \item Chose a \emph{suitable}\footnote{Matches the previous choice} probability distribution (must be a member of the EFDs). \dcheckmark
  \item Calculate the \emph{expected value} of this distribution for a particular
    data point. \dcheckmark
  \item Use \emph{maximum likelihood} to obtain an estimator of this quantity.
  \item Obtain an \emph{error function} as $E = -2\ln \mathcal{L}$ (deviance).
  \item Minimize the error $E$ by (typicall) Newton-Raphson.
    \begin{enumerate}
      \item initialize $\theta (0) := \boldsymbol 0, \; t \gets 0$
        \begin{align*}
          \text{loop until convergence:} \\
          \theta (t+1) = \theta(t) - \nabla E(\theta(t)) \cdot H_\theta^{-1} E(\theta(t)) \\
          t \gets t + 1
        \end{align*}
        \begin{align*}
          H_\theta E &= \begin{pmatrix}
            \frac{\partial^2 E}{\partial \theta_1^2} & \cdots & \frac{\partial^2 E}{\partial \theta_1 \partial \theta_k} \\
            \vdots & \ddots & \vdots \\
            \frac{\partial^2 E}{\partial \theta_k \partial \theta_1} & \cdots & \frac{\partial^2 E}{\partial \theta_k^2}
          \end{pmatrix} \\
          \nabla E &= \begin{pmatrix}
            \frac{\partial E}{\partial \theta_1} \\
            \vdots \\
            \frac{\partial E}{\partial \theta_k}
          \end{pmatrix}
        \end{align*}
    \end{enumerate}
\end{enumerate}

\begin{itemize}
  \item Typically, in the GLM setting, Newton-Raphson yields updating
    steps that require the solution of \emph{weighted least squares problems}.
    (called the IRLS family of algorithms: Iterated Reweighted Least Squares)
\end{itemize}
 % Teoria
%! TEX root = ../000-main.tex
% 2022-10-07
\chapter[SVMs]{Support Vector Machines}
\index{SVM}

\section{Linear SVM}
\label{sec:linear-svm}
\index{SVM!linear}

\begin{equation*}
	D = \{ (x^1,\, y^1),\, \dots,\, (x^n,\, y^n) \} \quad  x^i \in \mathds{R}^d \; y^i \in \{-1,\, 1\} \mid i = 1,\, \dots,\, n
\end{equation*}

The sum is a \iemph{binary classification}:
\begin{equation}
	\mathcal{F}_{\text{lin}} = \left\{
    x \to f\left(\omega^Tx + b\right) \mid
    \omega \in \mathds{R}^d,
    \; b \in \mathds{R}
  \right\}
\end{equation}

The equation $\pi : \omega^Tx + b = 0$ defines a hyperplane.
The idea is to predict class $+1$ whenever $\omega^Tx + b \geq 0$ and class $-1$ otherwise.

\begin{definition}{linsep}{}
	A dataset is linearly separable (\iemph{linsep}) if:
	\begin{equation*}
		\exists (\omega,\, \omega_0) \in \mathds{R}^d \setminus \{0\} \times \mathds{R}
		\text{ such that: }
		\forall i=1,\,\dots,\,n\quad y^i(\omega^Tx^i +b) > 0
	\end{equation*}
\end{definition}

We assume from now on that $D$ is \emph{linsep}.

\begin{definition}{geometric margin}{}
	The \iemph{geometric margin} of a linear
  classifier $f \to \omega^Tx + b$ \emph{at a data point $x$} is the
	euclidean distance from $x$ to the hyperplane $\pi$:
	\begin{equation*}
		\gamma(x) = \frac{\lVert\omega^Tx + b\rVert}{\lVert\omega\rVert} \tag{geometric margin at $x$}
	\end{equation*}

  \tcbline

  The \iemph{geometric margin of a linear classifier} on $D$ is the minimum of
these geometric margins over all data points. That is, it is the distance
from the closest point to the hyperplane:
\begin{equation*}
	\gamma = \min_{x \in D} \gamma(x) \tag{geometric margin of a linear classifier}
\end{equation*}

\end{definition}

\index{SVM!hyperplane}
\begin{figure}[H]
	\begin{tikzpicture}
		\begin{axis}[
				domain=0:11,
				xmin=0, xmax=10,
				ymin=0, ymax=10,
				samples=100,
				% axis y line=center,
				% axis x line=bottom,
				xlabel = {$n$},
				legend pos=south east,
				% area style,
				% ticks=none,
			]
			\addplot+[mark=none] {-x + 5};
			\addplot+[mark=none] {x};
			\addlegendentry{\(\pi: \omega^Tx + b = 0\)};
			\addlegendentry{\omega};
		\end{axis}

	\end{tikzpicture}
	\caption{Hyperplane and its normal vector}
\end{figure}

We set into the problem of finding the hyperplane $f$ that minimizes
its geometric margin over the data, $\gamma_{\text{min}}(f)$.

\begin{equation*}
	\max_{(\omega,\, b)}
	\left\{
	\min_{i=1,\,\dots,\, n}
	\frac{| \omega^Tx^i + b |}{\lVert\omega\rVert}
	\right\}
\end{equation*}

\subsection*{Questions}

\begin{enumerate}
	\item Why should we be doing this? ???
	\item Does a solution always exists? Yes
	\item Is the solution unique? Yes
	\item How do we find the solution?

	      The problem boils down to:
	      \begin{align*}
		      \max_{(\omega, b)}
		      \left\{
		      \min_{\substack{i=1,\dots, n                                                  \\ y^i(\omega^Tx^i + b) > 0}}
		      \underbrace{\frac{1}{\lVert\omega\rVert}}_{\text{margin of the solution}}
		      \right\}                                                                      \\
		      \min_{(\omega, b)} \frac{1}{2}\lVert\omega\rVert^2                            \\
		      \text{subject to} \quad y^i(\omega^Tx^i + b) \geq 1 \quad \forall i=1,\dots,n \\
		      \begin{rcases*}
			      \text{convex} \\
			      \text{QP}     \\
		      \end{rcases*} \Rightarrow \text{Solution is unique}
	      \end{align*}

	      \begin{enumerate}
		      \item Convex
		      \item Quadratic Programming (QP)
	      \end{enumerate}

\end{enumerate}

\subsubsection{The dual problem}
\index{SVM!dual}
\index{SVM!lagrangian function}
\index{SVM!lagrangian multiplier}
\index{SVM!KKT}

The Lagrangian function:
\begin{align*}
	\mathscr{L}(\omega, b, \alpha)                   & = \frac{1}{2}\lVert\omega\rVert^2 - \sum_{i=1}^n \alpha_i \left[ y^i(\omega^Tx^i + b) - 1 \right] \\
	\alpha                                           & = \begin{pmatrix} \alpha_1 \\ \vdots \\ \alpha_n \end{pmatrix}
	\tag{Lagrangian multiplier}                                                                                                                          \\
	\frac{\partial \mathscr{L}}{\partial \omega} = 0 & \Rightarrow \omega = \sum_{i=1}^n \alpha_i y^i x^i \tag{p1}\label{eq:lagp1}                       \\
	\frac{\partial \mathscr{L}}{\partial b} = 0      & \Rightarrow \sum_{i=1}^n \alpha_i y^i = 0 \tag{p2}\label{eq:lagp2}                                \\
	\text{KKT} \begin{pmatrix}
		           \text{Karush} \\
		           \text{Kuhn}   \\
		           \text{Tucker}
	           \end{pmatrix} \text{ conditions}                                                                                                          \\
	\forall i=1,\dots,n \quad \alpha_i \left[ y^i(\omega^Tx^i + b) - 1 \right] = 0 \tag{kkt}\label{eq:kkt}
\end{align*}

\begin{enumerate}
	\item \ref{eq:lagp1} \textrightarrow{} The solution $\omega$ must be a linear combination of the TR data points $x^i$
	\item \ref{eq:lagp2} \textrightarrow{} The sum of the Lagrangian multipliers is the same for both classes.
	\item \ref{eq:kkt} \textrightarrow{} $\alpha_i = 0$  xor $y^i(\omega^Tx^i + b) = 1$
	      \begin{enumerate}
		      \item $\alpha_i = 0$ \textrightarrow{} The data point $x^i$ is not a support vector
		      \item $y^i(\omega^Tx^i + b) = 1$ \textrightarrow{} The data point $x^i$ is a support vector. And coincide with the points on the margins of the hyperplanes.
	      \end{enumerate}
\end{enumerate}

% 2022-10-14

\begin{align*}
  \mathscr{L}_\text{dual}(\alpha) = \sum_{i=1}^n \alpha_i - \frac{1}{2} \mathlarger{\sum_{i=1}^n} \sum_{j=1}^n \alpha_i \alpha_j y^i y^j (x^i)^T x^j \\
	\text{goal} \quad \max_{\alpha} \mathscr{L}_\text{dual}(\alpha)                                                                       \\
	\text{subject to} \quad
	\begin{rcases*}
		\sum_{i=1}^n \alpha_i y^i = 0 \\
		\alpha_i \geq 0
	\end{rcases*} \quad \forall i=1,\dots,n
\end{align*}

\begin{note}{}{}
	The data points $x^i, i=1,\dots,n$  appear \emph{only} as
	dot products.
\end{note}

\begin{note}{}{}
	The (primal) solution vector $\omega$
	can be expressed as a linear combination of the (known) data points $x^i$.
\end{note}

% How \emph{general} are these properties?

\begin{question}{How \emph{general} are these properties?}{}
\end{question}

\begin{question}{What kind of problem do we have?}{}
	Quadratic problem subject to linear (inequality) constraints is
	a \iemph{quadratic programming} (QP) problem.
\end{question}

\index{SVM!primal-dual}
\begin{question}{What is the SVM model?}{}
	We have two versions of the SVM model, the primal and the dual:
	\begin{equation*}
		f_\text{SVM}(x) = \text{sign} \left( \hat{\omega}^T \boldsymbol{x} + \hat{b} \right)
		\quad \hat\omega \in \mathds{R}^d
		\quad b \in \mathds{R}
		\tag{PRIMAL}
	\end{equation*}

	For the \iemph{primal} version, we have $d+1$ parameters to estimate.

	\begin{align*}
		 & = \text{sign} \left( \left(\sum_{i=1}^n \hat\alpha_i y^i x^i \right)^T x + \hat{b} \right) \\
		 & = \text{sign} \left( \sum_{i=1}^n \hat\alpha_i y^i
		\underbrace{(x^i)^T x}_{\text{data}} + \hat{b} \right)
		\tag{DUAL}
	\end{align*}

	For the dual problem, we have $n$ parameters to estimate.
	And the data appears only as dot products.

	\vspace{0.5em}
	We use primal or dual depending on the dataset:
	\begin{enumerate}
		\item If $n \ll d$ we use the primal problem.
		\item If $n \gg d$ we use the dual problem.
	\end{enumerate}
\end{question}

\begin{question}{What kind of model is the dual version of $f_\text{SVM}$?}{}
	\begin{itemize}
		\item Linear (def.~\ref{def:linmod})
		\item Non-Parametric (def.~\ref{def:parametric})
		\item Discriminative (def.~\ref{def:discriminative})
	\end{itemize}
\end{question}

\begin{definition}{Parametric model}{parametric}
	A model is \iemph{parametric} when its size does not
	grow with the size of the training data.
	\begin{enumerate}
		\item size of the model = number of free parameters (unknowns)
		\item size of the training data = $n$
	\end{enumerate}
\end{definition}

\begin{definition}{Leave One Out cross-validation error}{loo}
	\index{LOO}
	$K$-fold cross-validation taken to the extreme: $K=n$.
	Means that we train on $n-1$ data points and test on the remaining one.
	We do this $n$ times, each time leaving out a different data point.
	The LOO error is the average of the $n$ test errors.
\end{definition}

\begin{theorem}{LOO analysis}{looanal}\index{LOO!analysis}

	\begin{lemma}{}{}
			The average LOO-CV error $f_n$ data sets of size $n \geq 2$
			is an unbiased estimator of the average (true) generalization
			error for data sets of size $n=1$.
		\begin{align*}
			\mathds{E}_{D^n \sim p^n}
			\left[ \hat{R}_{LOOCV}(f_{D_n}) \right] =
			\mathds{E}_{D^{n-1} \sim p^{n-1}} \left[{R_{LOOCV}(f_{D_n})} \right]
		\end{align*}
	\end{lemma}
	Let $f_{D_n}$ as in the lemma, let $SV(f_{D_n})$ be the set of
	support vectors in $D_n$. Then:
	\begin{equation*}
		\mathds{E}_{D^{n-1} \sim p^{n-1}} \left[{R(f_{D_n})} \right]
		\leq \mathds{E} \left[
			\frac{\left| SV(f_{D_n}) \right|}{n + 1}
			\right]
	\end{equation*}
  \tcblower
  \begin{note}
	  This is an argument in favor of \iemph{sparsity} (def.~\ref{def:sparsity}).
  \end{note}
\end{theorem}

\begin{definition}{Sparsity}{sparsity}\index{sparsity}
	The fraction of data points used in the final model.
\end{definition}

\subsection*{Non-separable data}

This means that there is at least one pair of data point $(x^i, y^i)$
such that $\nexists (\omega, b) \mid y^i (\omega^T x^i + b) \ngeq 1$.

The solution is:

\begin{equation*}
	y^i (\omega^T x^i + b) \geq 1 - \xi_i, \quad \xi_i \geq 0,\quad
	i=1,\dots,n
\end{equation*}

These $\xi_i$ variables are called \iemph{slack variables}.

$\xi_i$ is the amount needed s.t. $y^i (\omega^T x^i + b) \geq 1$.

\begin{figure}[H]
	\begin{tikzpicture}[
			important line/.style={thick,blue,thick},
			dashed line/.style={dashed,blue,thick},
			leftNode/.style={circle,minimum width=.5ex, fill=blue,draw},
			rightNode/.style={rectangle,minimum width=.5ex, fill=red,thick,draw},
		]
		\pgfmathsetseed{1234}
		\begin{axis}[
				width=0.8\textwidth,
				xmin=-50,xmax=50,
				ymin=-50,ymax=50,
				ticks=none,
				axis lines=none,
			]
			\pgfmathsetmacro{\sep}{20};
			\pgfmathsetmacro{\xorig}{120};
			\pgfmathsetmacro{\rot}{60};

			\begin{scope}[rotate around={\rot:(0,0)}]
				\draw[dashed line] (-100, -\sep) -- (100, -\sep);
				\draw[dashed line] (-100, \sep) -- (100, \sep);
				\draw[important line] (-100, 0) -- (100, 0);

				\node[leftNode,label={$x_1$},name=x1] at (-20, \sep) {};
				\node[rightNode,label={$x_2$},name=x2] at (-22, -\sep) {};

				\node[leftNode,label={$x_3$},name=x3] at (-10, -\sep*3/2) {};
				\node[leftNode,label={$x_5$},name=x5] at (-27, -\sep/3) {};

				\node[rightNode,label={$x_4$},name=x4] at (10, -\sep*2/3) {};

				\coordinate (origin) at (0, 0);
				\coordinate (left) at (0, \sep);
				\coordinate (right) at (0, -\sep);

				\coordinate (x3l) at (x3 |- left);
				\coordinate (x4r) at (x4 |- right);
				\coordinate (x5l) at (x5 |- left);

				\begin{scope}[color=purple]
					\draw (x3) -- (x3l) {};
					\draw (x4) -- (x4r) {};
					\draw (x5) -- (x5l) {};
					\node[anchor=north] at ($(x3)!0.5!(x3l)$) {$\xi_3$};
					\node[anchor=north] at ($(x4)!0.5!(x4r)$) {$\xi_4$};
					\node[anchor=north] at ($(x5)!0.5!(x5l)$) {$\xi_5$};
				\end{scope}

				\draw[<->,orange,line width=2pt] (-34,\sep) -- (-34,-\sep);
				\node[orange!90!black,anchor=north,rotate=\rot-90,fill=white] at (-34.25,0) {margin = $\frac{2}{\lVert \omega \rVert}$};

				\pgfplotsinvokeforeach{0.00,0.1,...,1.00}{
					\node [leftNode] at (rand*30-20,\sep+rnd*30) {};
					\node [rightNode] at (rand*30,-\sep-rnd*30) {};
				}

			\end{scope}
			\node[blue,anchor=south east,rotate=\rot] at (2,45) {$\omega^T x + b = 1$};
			\node[blue,anchor=south east,rotate=\rot] at (22,45) {$\pi:\,\omega^T x + b = 0$};
			\node[blue,anchor=south east,rotate=\rot] at (42,45) {$\omega^T x + b = -1$};

		\end{axis}
	\end{tikzpicture}
	\caption{SVM Slack Variables}\label{fig:slack}
\end{figure}

If we remove the points inside the margin, we get a new dataset
which is linearly separable and has 0 training error. This is called
\iemph{soft-margin SVM}, in contrast with a hard-margin SVM, and
the margin will be $\frac{2}{\|\omega\|}$.

Finding the hyperplane with the smallest training error
is NP-hard.

We  have two antagonic goals:
\begin{enumerate}
	\item maximize the margin
	\item \textbf{minimize the training error} An upper bound on the
	      number of training errors is upper bounded by the sum of the slacks:
	      \begin{equation*}
		      n TR \leq \sum_{i=1}^n \xi_i
	      \end{equation*}
\end{enumerate}

The new optimization problem is:
\begin{align*}
	\min_{\omega, b, \xi} &
	\underbrace{\frac{1}{2} \|\omega\|^2}_{\substack{\text{maximize}     \\\text{margin}}}
	+ C \cdot \underbrace{\sum_{i=1}^n \xi_i}_{\substack{\text{minimize} \\\text{TR}\\\text{upper bound of TR}}} \\
	\text{s.t.}           & \begin{rcases*}
		                        y^i (\omega^T x^i + b) \geq 1 - \xi_i \\
		                        \xi_i \geq 0
	                        \end{rcases*}
	\quad i=1,\dots,n
\end{align*}

\index{SVM!trade-off}
We add a hyperparameter $C$ which controls the trade-off between
these two goals.

% 2022-10-21
\subsection{Loss functions for the SVM (classification case)}
\index{SVM!loss functions}

\begin{align*}
	\lVert \xi \rVert_q^q          & = \sum_{i=1}^n \xi_i,\quad q \geq 1      \\
	\lVert \xi \rVert_\infty       & = \max \{ \xi_1,\,\xi_2,\,\dots \xi_n \} \\
	\text{Typical values for $q$:} & \begin{cases*}
		                                 q=1 & \text{1-norm SVM} \\
		                                 q=2 & \text{2-norm SVM} \\
	                                 \end{cases*}
\end{align*}

\begin{figure}[H]
\begin{tikzpicture}
	\pgfmathsetseed{1234}
	\begin{axis}[
			width=0.8\textwidth,
      height=0.5\textwidth,
      axis equal,
			xmin=-3,xmax=3.5,
			ymin=-1,ymax=5,
			xtick={0,1},
			ytick={0,1},
			axis lines=middle,
			xlabel=$z$,
      yticklabel style={below left},
      legend pos=north east,
      legend cell align={left}
		]
		% quadratic hinge loss
		\draw[green!70!black, line width=2pt] (1,0) -- (2,0);
		\addplot[green!70!black, line width=2pt, domain=-3:1] {(1 - x)^2};
		\addlegendentry{quadratic hinge loss}

		% hinge loss
		\draw[blue, line width=2pt] (-4,5) -- (1,0) -- (2,0);
		\addlegendimage{draw=blue, line width=2pt};
		\addlegendentry{hinge loss}

		% 0 1 loss
		\draw[red, line width=2pt] (-5,1) -- (0,1);
		\draw[red, line width=2pt](0,0) -- (2,0);
		\draw[red, line width=2pt,dotted](2,0) -- (3,0);
		\addlegendimage{draw=red, line width=2pt};
		\addlegendentry{0-1 loss}


		\node[red, fill=white, draw, circle, inner sep=0pt, minimum size=5pt] at (0,1) {};
		\node[red, fill=red, draw, circle, inner sep=0pt, minimum size=5pt] at (0,0) {};
	\end{axis}
\end{tikzpicture}
\caption{Loss functions for the SVM (classification case)}\label{fig:loss}
\end{figure}

\begin{align*}
	L_{0/1} (z)  & = \begin{cases}
		                 0 & \text{if } z \geq 0 \\
		                 1 & \text{if } z < 0
                   \end{cases} \equiv \mathds{1}_{\{z \leq 0\}}          \\
	L_{hl} (z)   & = \max \{ 0, 1 - z \}                &  & (q=1) \\
	L_{hl^2} (z) & = \max \{ 0, 1 - z \}^2              &  & (q=2) \\
\end{align*}

\begin{exercise}{Homework}{}
	repeat all the steps from the primal to the final solution
	(model): exactly the same!

	\begin{equation*}
		0 \leq \alpha_i \leq C
	\end{equation*}
\end{exercise}

\pagebreak
\subsection{SVMs for regression}
\index{SVMR}
\index{SVM!regression}
\begin{marker}
\paragraph*{Idea}: we want a regression algorithm that maintains the defining
(and good) features of the SVMC\footnote{Support Vector Machine for Classification}:
\end{marker}
\begin{itemize}
	\item \textbf{Sparsity} and obtention of SVs (Support Vectors).
	\item \textbf{Notion of margin} as a regularizer.
	\item \textbf{Unique solution} (QP)
	\item Data points $x_i$ appear in the model ($f_{SVM}$) and the QP problem
	      \emph{only} as linear products.
	\item Linear models
\end{itemize}

\begin{figure}[H]
	\begin{tikzpicture}[
			important line/.style={thick,blue,thick},
			dashed line/.style={dashed,blue,thick},
			leftNode/.style={circle,minimum width=.5ex, fill=blue,draw},
			rightNode/.style={rectangle,minimum width=.5ex, fill=red,thick,draw},
		]
		\pgfmathsetseed{1234}
		\begin{axis}[
				width=0.8\textwidth,
				xmin=-50,xmax=50,
				ymin=-40,ymax=40,
				ticks=none,
				axis lines=none,
        legend pos=north west,
			]
			\pgfmathsetmacro{\rot}{45};

			\begin{scope}[rotate around={\rot:(0,0)}]

				\draw[fill=blue,opacity=0.4] (-80,-10) rectangle (80,10);

        \draw[line width=2pt] (-80,0) -- (80,0);
				\draw[dashed, line width=2pt] (-80,10) -- (80,10);
				\draw[dashed, line width=2pt] (-80,-10) -- (80,-10);

				\pgfplotsinvokeforeach{1,...,30}{
					\node [rightNode] at (rand*50, rand*15) {};
				};
				\pgfplotsinvokeforeach{1,...,5}{
					\node [rightNode] at (rand*50, rnd*30 + 10) {};
					\node [rightNode] at (rand*50, -rnd*30 - 10) {};
				};

      % \addplot[only marks,red,mark=*,mark size=1pt,domain=-50:50] {rand*30};

        \addlegendimage{legend image code/.code={\node [rightNode,anchor=center] at (2mm,0){};}};
        \addlegendentry{data};
				\addlegendimage{area legend,fill=blue,opacity=0.4};
				\addlegendentry{$\varepsilon$-tube};

			\end{scope}

		\end{axis}
	\end{tikzpicture}
  \caption{SVMR: regression with an $\varepsilon$-tube}
\end{figure}

\begin{itemize}
	\item $\varepsilon > 0$
	\item All points strictly inside the $\epsilon$-tube entail 0 error.
	\item All points (strictly) outside the $\epsilon$-tube entail a positive,
	      equal to the vertical distance to the nearest margin side.
\end{itemize}

\begin{equation*}
	L_{\epsilon} (y,\,y') = \begin{cases}
		0                      & \text{if } |y - y'| \leq \varepsilon \\
		|y - y'| - \varepsilon & \text{otherwise }
	\end{cases}
\end{equation*}

We need to introduce two sets of slacks,
$\overline{\xi_i},\, \underline{\xi_i} \geq 0,\; i=1,\dots,n$,
\index{SVM!slacks}

\begin{align*}
	\min_{w,\,b,\,\xi,\,\overline{\xi},\,\underline{\xi}} &
	\left\{ \frac{1}{2} \|w\|^2 +
	C \sum_{i=1}^n \left(\overline{\xi}_i + \underline{\xi}_i \right) \right\}                                                  \\
	\text{s.t.}                                           & \quad \begin{rcases*}
		                                                              (\omega^T x_i + b) \geq y_i + \varepsilon + \overline{\xi}_i  \\
		                                                              (\omega^T x_i + b) \leq y_i - \varepsilon - \underline{\xi}_i \\
		                                                              \overline{\xi}_i,\, \underline{\xi}_i \geq 0
	                                                              \end{rcases*} \quad i=1,\dots,n \\
                                                        && \text{Convex QP problem}
\end{align*}

\begin{align*}
	f_{SVMR} (x) & = \sum_{i=1}^n (\overline{\alpha}_i - \underline{\alpha}_i)
	(x^i)^T x + b =                                                                                                                           \\
	             & \downarrow \beta_i = \overline{\alpha}_i - \underline{\alpha}_i                                                            \\
	             & = \sum_{i=1}^n (\beta_i) (x^i)^T x + b                                                                                     \\
	             &                                                                 & \text{Many of }\beta_i\text{ are zero}                   \\
	             &                                                                 & \text{SVs are those }x^i/\beta_i \neq 0                  \\
	             &                                                                 & \text{We have two hyper-parameters: }C,\,\varepsilon > 0 \\
	             &                                                                 & \text{overfitting } C \to \infty,\, \varepsilon \to 0
\end{align*}


%! TEX root = ../000-main.tex
\section{Model complexity, VC-dimension and generalization bounds}
\index{VC-dimension}\index{complexity}\index{generalization bound}

\begin{definition}{dichotomy of a set}{dichotomy}

  Let $\mathcal{F}$ be a hypothesis set, let $S$ be a set of data points $S =
  \{x_1, \ldots, x_m\},\quad x^i \in \mathds{R}^d$.

  A \iemph{dichotomy} of a set
  $S$ is one of the possible ways of labelling the points in $S$ in
  one of two classes.

\end{definition}

\begin{definition}{Shattering}{shattering}\index{shattering}
  A set $S$ of size $m$ is said to \emph{shatter} by $\mathcal{F}$
  when all $2^m$ possible dichotomies can be realized (makes 0 TR on $S$)
  by a member of
  $\mathcal{F}$.
\end{definition}

\begin{definition}{VC-dimension}{vc-dimension-cont}
  Vapnik-Chervonenkis dimension \\[1em]
  The VC-dimension of a hypothesis space $\mathcal{F}$ is the size
  of the largest set $S$ that can be shattered by $\mathcal{F}$.
\end{definition}

\begin{ex}
  Intervals on the real line ($d = 1$)
  \begin{equation*}
    \mathcal{F} = \{[a,\,b]\mid a < b \in \mathds{R}\}
  \end{equation*}

  The VC-dimension of $\mathcal{F}$ is $2$. With $m = 3$ points,
  we cannot shatter it. But with $m = 2$ we can:
\begin{figure}[H]
  \begin{tikzpicture}[
			leftNode/.style={circle,minimum width=.5ex, fill=red,draw},
			rightNode/.style={circle,minimum width=.5ex, red, fill=white,thick,draw},
      limitNodeL/.style={draw,minimum width=0.25ex,minimum height=1ex,fill,label=$a$},
      limitNodeR/.style={draw,minimum width=0.25ex,minimum height=1ex,fill,label=$b$},
    ]

    % \addlegendimage{leftNode};
    % \addlegendentry{+1};
    % \addlegendimage{rightNode};
    % \addlegendentry{-1};

    \draw[->] (-1,0) -- (5,0);
    \draw[->] (-1,1) -- (5,1);
    \draw[->] (-1,2) -- (5,2);

    \node[limitNodeL] at (0,0) {};
    \node[limitNodeR] at (4,0) {};
    \draw[line width=2pt] (0,0) -- (4,0);
    \node[leftNode] at (1,0) {};
    \node[leftNode] at (3,0) {};

    \node[limitNodeL] at (0,1) {};
    \node[limitNodeR] at (2,1) {};
    \draw[line width=2pt] (0,1) -- (2,1);
    \node[leftNode] at (1,1) {};
    \node[rightNode] at (3,1) {};

    \node[limitNodeL] at (-0.5,2) {};
    \node[limitNodeR] at (0.5,2) {};
    \draw[line width=2pt] (-0.5,2) -- (0.5,2);
    \node[rightNode] at (1,2) {};
    \node[rightNode] at (3,2) {};

    \matrix [draw,right] at (6,1) {
      \node [leftNode,label=right:{$+1$}] {}; \\
      \node [rightNode,label=right:{$-1$}] {}; \\
    };

	\end{tikzpicture}
\end{figure}
\end{ex}

\begin{ex}
  Axis aligned rectangles in $\mathds{R}^2$ ($d = 2$)
  \begin{equation*}
    \mathcal{F} = \{[a_1,\,b_1]\times [a_2,\,b_2]\mid a_1 < b_1,\, a_2 < b_2 \in \mathds{R}\}
  \end{equation*}

  The VC-dimension of $\mathcal{F}$ is $4$.

\begin{figure}[H]
  \begin{tikzpicture}[
			nodeIn/.style={circle,minimum width=.5ex, fill=red,draw},
			nodeOut/.style={circle,minimum width=.5ex, red, fill=white,thick,draw},
      limitNodeL/.style={draw,minimum width=0.25ex,minimum height=1ex,fill,label=$a$},
      limitNodeR/.style={draw,minimum width=0.25ex,minimum height=1ex,fill,label=$b$},
    ]

    \node[nodeIn] at (0, 1) {};
    \node[nodeIn] at (0, -1) {};
    \node[nodeOut] at (1, 0) {};
    \node[nodeOut] at (-1, 0) {};
    \draw (-0.5, -1.5) rectangle (0.5, 1.5);
      \draw[<->] (-1.75, 1.75) -- (-1.75, -1.75) -- (1.75, -1.75);

    \begin{scope}[shift={(-4,0)}]
      \node[nodeIn] at (0, 1) {};
      \node[nodeIn] at (0, -1) {};
      \node[nodeOut] at (1, 0) {};
      \node[nodeIn] at (-1, 0) {};
    \draw (-1.5, -1.5) rectangle (0.5, 1.5);
      \draw[<->] (-1.75, 1.75) -- (-1.75, -1.75) -- (1.75, -1.75);
    \end{scope}

    \begin{scope}[shift={(-4,-4)}]
      \node[nodeIn] at (0, 1) {};
      \node[nodeIn] at (0, -1) {};
      \node[nodeIn] at (1, 0) {};
      \node[nodeIn] at (-1, 0) {};
    \draw (-1.5, -1.5) rectangle (1.5, 1.5);
      \draw[<->] (-1.75, 1.75) -- (-1.75, -1.75) -- (1.75, -1.75);
    \end{scope}

    \begin{scope}[shift={(0,-4)}]
      \node[nodeIn] at (0, 1) {};
    \draw (-0.5, 0.5) rectangle (0.5, 1.5);
      \node[nodeOut] at (0, -1) {};
      \node[nodeOut] at (1, 0) {};
      \node[nodeOut] at (-1, 0) {};
      \draw[<->] (-1.75, 1.75) -- (-1.75, -1.75) -- (1.75, -1.75);
    \end{scope}

    \matrix [draw,right] at (2,-2) {
      \node [nodeIn,label=right:{$+1$}] {}; \\
      \node [nodeOut,label=right:{$-1$}] {}; \\
    };

	\end{tikzpicture}
  \caption{Axis aligned rectangles in $\mathds{R}^2$}
\end{figure}
\end{ex}

\begin{note}
  Radamaiher complexity encompases the VC-dimension adding
  regression.
\end{note}

%! TEX root = ../000-main.tex
\chapter[RKHSs]{Learning in RKHSs}

\section{Introduction}

\subsection{Hilbert Spaces}

\begin{definition}[parbox=false]{Hilbert space (HS)}{hilber-space}
	\begin{itemize}
		\item Consider a space $G$ of functions $f: \mathcal{X} \to \mathbb{R}$,
		      where $\mathcal{X}$ is the data space.

		\item Define an \iemph{inner product} (ip) $\langle f, g \rangle, \; f, g \in G$ % by ???
		\item $\{f_n\}_{n \in \mathds{N}}$ converges to $f$ iff
		      $$\forall \epsilon > 0, \exists N \in \mathds{N} \text{ s.t. } \forall n \geq N, \; \lVert f_n - f \rVert < \epsilon$$
	\end{itemize}
	\begin{definition}{Cauchy sequence}{cauchy-sequence}\index{Cauchy sequence}
		$\{f_n\}_{n \in \mathds{N}}$ is Cauchy when:

		$$\forall \epsilon > 0, \; \exists N \in \mathds{N} \text{ s.t. } \forall n, m \geq N, \; \lVert f_n - f_m \rVert < \epsilon$$
	\end{definition}

	A space in which all Cauchy sequences converge to an element of the space is
	called \emph{complete}.\index{complete space}

	A complete vector space in which we have
	a norm ($\lVert \cdot \rVert$) is called a \iemph{Banach space}.

	\begin{marker}
		A \iemph{Hilbert space} (HS) is a Banach space where the norm
		comes from an inner product.
	\end{marker}

\end{definition}

\begin{question}{Homework}{}
	Is there an implication between convergence and Cauchy convergence?
\end{question}

\begin{example}{Hilbert spaces}{}
	\begin{itemize}
		\item $\mathds{R}^d, \; d \in \mathds{N}$
		\item The space of square-integrable functions
	\end{itemize}
\end{example}

\subsection{Reproducing Kernel}

\begin{definition}[parbox=false]{RKHSs}{RKHSs}

	$\mathcal{X}$ data space ($\mathds{R}^d$)

	$\mathcal{H}$ of functions $f: \mathcal{X} \to \mathds{R}$
	\tcbline
	A \iemph{Reproducing Kernel Hilbert space} (RKHS) is a Hilbert space
	of functions in which all linear evaluation functionals are continuous:
	\begin{align*}
		[x](\cdot) : \mathcal{H} & \to \mathds{R}                                     \\
		f                        & \mapsto [x](f) = f(x)                              \\
		\span
		\text{ is continuous } \forall f \in \mathcal{H} \; \forall x \in \mathcal{X} \\
		\iff \text{ all functionals are bounded }
	\end{align*}
\end{definition}

\begin{theorem}{}{}
	In a RKHS $\mathcal{H}$, $\forall x \in \mathcal{X}$, there
	exists an element $\kappa_x(\cdot) \in \mathcal{H}$, which
	``represents'' the evaluation functional $[x](\cdot)$
	in this sense:
	\begin{align*}
		\langle \kappa_x(\cdot), f \rangle_{\mathcal{H}} = f(x) \forall f \in \mathcal{H}
	\end{align*}
	\begin{marker}
		$\kappa_x(\cdot)$ is a function, we add the $(\cdot)$ to emphasize it.
	\end{marker}
\end{theorem}


\begin{definition}{Reproducing kernel}{}
	The symmetric function
	\begin{equation*}
		\kappa(x, x') \coloneqq \kappa_x(x'),\; \forall x, x' \in \mathcal{X}
	\end{equation*}
	is called the \iemph{reproducing kernel} of the space (rk).
\end{definition}

\subsection{Positive Semi-Definite Functions}

\begin{definition}{Positive semi-definite functions}{psd}
	Consider a data space $\mathcal{X}$. A symmetric bivariate function
	$\kappa: \mathcal{X} \times \mathcal{X} \to \mathds{R}$ is
	\iemph{positive semi-definite} (PSD) when
	$\forall n \in \mathds{N}^+, \forall \{x^1,\, x^2,\,\dots, x^n\},
		x^i \in \mathcal{X}$.
	\begin{align*}
		\mathlarger{\sum_{i=1}^n} \sum_{j=1}^n  a_i a_j \kappa(x^i, x^j) \geq 0,\quad
		\forall a=\begin{pmatrix} a_1 \\ \vdots \\ a_n \end{pmatrix},\; a_i \in \mathds{R}
	\end{align*}

	\tcblower

	(alternative) If we form the following matrix:
	\begin{align*}
		\mathlarger{K}_{ij} \coloneqq \kappa(x^i, x^j), \\ \text{ then } a^TKa \geq 0
	\end{align*}
\end{definition}

\begin{example}{}{}
	\begin{align*}
		\mathcal{X} & = \{1, 2, 3, \dots, N \} \rightarrow{}
		\text{ In this case } \kappa \text{ is empty } K     \\
		\mathcal{X} & = [a,b] \subset \mathds{R}             \\
		\mathcal{X} & = \mathds{R}^d
	\end{align*}
\end{example}

\begin{theorem}[parbox=false]{Moore-Aronszajn '50s}{}
	For every PSD funciton $\kappa: \mathcal{X} \times \mathcal{X} \to \mathds{R}$, there
	is a \emph{unique} RKHS $\mathcal{H}_\kappa$ for which $\kappa$ is the reproducing kernel.

	\tcblower

	More precisely,

	Define $\kappa_x(\cdot) \coloneqq \kappa(x, \cdot)$; then
	$\forall f \in \mathcal{H}_\kappa,
		\langle \kappa_x, f \rangle = f(x), \forall x \in \mathcal{X}$.

	\begin{note}The converse is also true:
		If we depart from a RKHS $\mathcal{H}_{\kappa}$ having
		$\kappa$ as reproducing kernel, then $\kappa$ is PSD function.
	\end{note}
\end{theorem}

\pagebreak
\section{Construction of the RKHS}

Depart from a PSD function $\kappa: \mathcal{X} \times \mathcal{X} \to \mathds{R}$.
And a dataset $\mathcal{D} = \{x^1, \dots, x^n\}$.

\begin{itemize}
	\item $\mathcal{H}_\kappa \coloneqq \emptyset$.
	\item $\forall x \in \mathcal{X}$, define $\kappa_x(\cdot) \coloneqq \kappa(x, \cdot)$.
	      and $\forall x \in \mathcal{X}$, add $\kappa_x$ to $\mathcal{H}_\kappa$.
	\item Add to $\mathcal{H}_\kappa$ its own \iemph{span}:
	      \begin{equation*}
		      \mathcal{H}_\kappa \coloneqq \mathcal{H}_\kappa \cup \{
		      f(\cdot) = \sum_{i=1}^n a_i \kappa_{x^i}(\cdot) \mid
		      \forall x \in \mathcal{X},
		      \forall n \in \mathds{N}^+,
		      \forall a_i \in \mathds{R},
		      i=1,\dots,n
		      \}
	      \end{equation*}
	\item We now define the inner product:
	      \begin{align*}
		      \langle f, g \rangle_{\mathcal{H}_\kappa} & \coloneqq
		      \left\langle
		      \sum_{i=1}^n a_i \kappa_{x^i}(\cdot), \sum_{j=1}^{n'} a'_j \kappa_{x'^j}(\cdot)
		      \right\rangle                                                                                                          \\
		                                                & = \mathlarger{\sum_{i=1}^n} \sum_{j=1}^{n'} a_i a'_j \kappa({x^i}, {x'^j})
	      \end{align*}
	\item Let's show that $\kappa$ is the r.k of $\mathcal{H}_\kappa$
	      \begin{align*}
		      \forall x                                        & \in \mathcal{X}                          \\
		      \langle \kappa_x, f \rangle_{\mathcal{H}_\kappa} & =
		      \left\langle \kappa_x, \sum_{i=1}^n a_i \kappa_{x^i} \right\rangle_{\mathcal{H}_\kappa}     \\
		                                                       & = \sum_{i=1}^n a_i \kappa(x, x^i) = f(x)
	      \end{align*}
	\item \iemph{Completeness}: We take all limits of Cauchy sequences
	      and add them to the space.
\end{itemize}

\begin{definition}{Span of a set}{}\index{Span of a set}
	The set of all linear combinations of the elements of the set.
\end{definition}

% 2022-11-18

\begin{theorem}{Representer theorem}{representer}\index{Representer theorem}
	(original by Kimmerldorf \& Wahba '70s)

	Let $\mathcal X$ be a data space and
	$\kappa : \mathcal X \times \mathcal X \to \mathds R$ a kernel function.

	We are given a training set $D_n =\{(x^1, y^1), \dots, (x^n, y^n)\},\,
		x^i \in \mathcal X,\, y^i \in \mathds R$.
	call $\mathcal H_\kappa$ the \emph{RKHS} generated by $\kappa$.
	let $\Omega: \mathcal H_\kappa \to \mathds R$ be a \emph{non-decreasing}
	function (the \emph{regularizer}), and
	$l: (\mathds R)^{2n} \to \mathds R^+_0$ be a \emph{loss function}.

	Then, the following \iemph{optimization problem}:
	\begin{align*}
		\argmin_{f \in \mathcal H_\kappa}
		\left\{
		l(f(x^1), y^1 , \ldots , l(f(x^n), y^n) + \Omega(\lVert f \rVert_{\mathcal H_\kappa})
		\right\}
	\end{align*}
	has solutions\footnote{There is always at least one solution with this form. (There could be more than one and others that do not have this form)} of the form:
	\begin{align*}
		f    & = \sum_{i=1}^n \alpha_i \kappa_{x^i} \in \mathcal H_\kappa \\
		f(x) & = \sum_{i=1}^n \alpha_i
		\underbrace{\kappa(x^i,\, x)}_{\kappa_{x^i}(x)}
	\end{align*}

	Moreover, if $\Omega$ is \emph{increasing}, then all solutions have
	this form.

	\begin{note}
		Typically, the optimization problem is not that general and takes
		the form:
		\begin{equation*}
			\argmin_{f \in \mathcal H_\kappa} \left\{
			\underbrace{\sum_{i=1}^n l(f(x^i), y^i)}_{\text{Minimize TR}} +
			\underbrace{\lambda \lVert f \rVert_{q\,\mathcal H_\kappa}^q}_{\text{Penalize complexity}}
			\right\}
		\end{equation*}
		where $\lambda$ is the \emph{regularization constant}.

		\tcbline

		In this setting, if the loss function $l$ is \emph{convex}
		with respect to the prediction $f(x)$, then the solution is
		\emph{unique}.
	\end{note}

\end{theorem}

\pagebreak
\begin{example}[breakable]{Kernel ridge Regression (with basis functions)}{}
	\tcbsubtitle{Standard Ridge Regression}
	\begin{align*}
		\min_{\omega \in \mathds R^D} \left\{
		\sum_{i=1}^n (f(x_i) - y_i)^2 + \lambda \lVert \omega \rVert^2
		\right\},\, \lambda > 0
	\end{align*}
	where:
	\begin{align*}
		f(x)      & = \omega^T \phi(x)                                    \\
		\phi :    & \mathds R^D \to \mathds R^D                           \\
		x \mapsto & \phi(x) = \begin{pmatrix}
			                      \phi_1(x) \\
			                      \vdots    \\
			                      \phi_D(x)
		                      \end{pmatrix},\, \phi \text{ is predefined}
	\end{align*}
	\tcbline
	In matrix form: We depart from:
	$X_{x\times d} \leadsto_{\phi} \Phi_{x\times D}$,
	such as: $\Phi_{ij} = \phi_j(x^i)$. Then
	% $f(x) = \omega^T \phi(x) = \omega^T \Phi^T x$.
	the solution is:
	\begin{equation*}
		\hat \omega \coloneqq
		(\underbrace{\vphantom{\bigr|}\Phi^T \Phi + \lambda I_D}_{\substack{\text{Positive definite}
				\\ \implies \text{non-singular}}})^{-1} \Phi^T y
	\end{equation*}

	\paragraph{Computational complexity}
	We have to solve a linear system of $D$ equations and $D$ unknowns,
	which is $O(D^3)$.

	\begin{note}
		This is the \emph{PRIMAL} formulation of ridge regression.
	\end{note}

	\tcbsubtitle{DUAL formulation}

	\begin{note}
		The key idea is that $\omega$ can be expressed as linear combination of
		the (transformed) data points:
		\begin{equation*}
			\omega = \sum_{i=1}^n \alpha_i \phi(x^i)
		\end{equation*}
		In matrix vector form:
		\begin{equation*}
			\omega = \Phi^T \alpha
		\end{equation*}
	\end{note}

	$\frac{\partial J}{\partial \omega} = 0$ leads to:
	\begin{equation*}
		\Phi^T\Phi \omega + \lambda \omega = \Phi^T y \tag{normal equations}
	\end{equation*}
	\begin{align*}
		\leadsto \omega = \frac{1}{\lambda} \Phi^T (y - \Phi \omega)
		\equiv \Phi \alpha                              \\
		\implies \lambda \alpha = y - \Phi\Phi^T \alpha \\
		\hat \alpha = (\Phi\Phi^T + \lambda I)^{-1}y
	\end{align*}

	\begin{note}
		$\Phi\Phi^T$ is a $n\times n$ matrix. In contrast, $\Phi^T\Phi$ is a
		$D\times D$ matrix.

		This is the \emph{DUAL} formulation of ridge regression.
		and is more efficient than the primal formulation. Since
		the cost is $O(n^3)$ instead of $O(D^3)$.
	\end{note}

	Left $G = \Phi\Phi^T$, then $G_{ij} = \phi(x^i)^T \phi(x^j)$.
	It is a matrix of \emph{inner products} between known
	data points. This is the \iemph{Gram matrix} (or the Gramian)

	\tcbline

	Now, we can write the solution in terms of the Gram matrix:
	\begin{align*}
		f(x) = \hat \omega^T \phi(x) & = (\Phi^T\hat\alpha)^T \phi(x) \\
		                             & =
		\underbrace{(\Phi^T(\Phi\Phi^T + \lambda I)^{-1}y)^T}_A \phi(x) = A\phi(x)
	\end{align*}
	\begin{align*}
		f(x) = \hat\omega^T \phi(x) & =
		\left(
		\sum_{i=1}^n \alpha_i \phi(x^i)
		\right)^T \phi(x)                                                             \\
		                            & = \sum_{i=1}^n \hat\alpha_i \phi(x^i)^T \phi(x) \\
	\end{align*}
	These are the inner products between data points.

	Now we \emph{kernelize} the inner products:

	Replace $\phi(x^i)^T \phi(x)$ with $\kappa(x^i, x)$, where $\kappa$ is a
	kernel function.

	The model becomes:

	\begin{equation*}
		f(x) = \sum_{i=1}^n \hat\alpha_i \kappa(x^i, x)
	\end{equation*}

	Now we are doing ridge regression not in $\mathds R^d$, but in
	$\mathds R^D$. But we are \emph{not} explicitly computing
	any $\phi(\cdot)$.

	This is what we call the \iemph{kernel trick}.

  Our intuition\footnote{VC-dimension should be kept low.
    The VC-dimension of a linear classifier is $d+1$
  }
  says that $D \gg d$. Actually, sometimes
	$D \in \mathds N \vee \{\infty\}$.
	\begin{align*}
		\begin{pmatrix}
			d + q \\
			q
		\end{pmatrix} & = D \in \mathds N \tag{Polynomial kernel} \\
		D               & = \infty \tag{RBF kernel}
	\end{align*}

  We are doing standard ridge regression in the Hilbert space
  $\mathcal H = \mathds R^D$, instead of the original space $\mathds R^d$.

  \begin{note}
    At no point we need to know that there is an inner product in the
    data space $\mathcal X$.
  \end{note}
\end{example}

\begin{example}{Same solution via the representer theorem}{}

  Let $\phi: \mathds R^d \to \mathds R^D$ be a mapping between
  the data space $\mathcal X$ and a Hilbert space $\mathcal H$.

  Let $\kappa : \mathcal X \times \mathcal X \to \mathds R$ be a kernel
  function.

  Define:
  \begin{equation*}
    \kappa(x, x') \coloneqq \langle \phi(x), \phi(x') \rangle_{\mathcal H_\kappa}
  \end{equation*}
  where $\mathcal H_\kappa$ is the \emph{RKHS} associated with $\kappa$.
  \begin{note}
    We use the abstract inner product $\langle \cdot, \cdot \rangle$ to
    denote the inner product in $\mathcal H$. (Not $\phi(x)^T \phi(x')$)
  \end{note}

  We have a data set $D_n$ and $l(f(x^i), y^i) = (f(x^i) - y^i)^2$.
  \begin{equation*}
    \Omega(\lVert f \rVert) =\lambda\lVert f \rVert_{\mathcal H_\kappa}^2
  \end{equation*}


\end{example}




\part{Problems}
% Problem sessions
\chapter{GLMs}

\section{Ridge Regression}

The optimal value for the regularization parameter $\lambda$ in \emph{ridge regression} is:

\begin{equation}
	\lambda^* = \frac{\sigma^2}{n \sigma_{\omega}^2}
\end{equation}

\subsection{Predictive distribution} % (fold)

\paragraph{``Full Bayessian''} Suppose we are given a \emph{new data point}
\(x^0 \to y^0\).

\begin{align*}
	P(y^0 \mid x^0, y, \alpha, \beta) & = \int_{\mathds{R}^{d+1}}
	\underbrace{P(y^0 \mid x^0, \underbar{\omega}, \beta)}_{\substack{\text{Prob. setting}        \\ \text{of my model}}}
	\cdot
	P(\underbar{\omega} \mid y, \alpha, \beta)
	d\underbar{\omega}                                                                         \\
	                               & = \mathcal{N}\left(y^0\mu^Tx^0,\; \sigma^2(x^0) \right)   \\
	\sigma^2(x)                    & = \underbrace{\frac{1}{\beta}}_\text{Noise in the data} +
	\underbrace{X^T\Sigma X \vphantom{\frac{1}{\beta}}
	}_\text{Uncertainty on the parameters $\underbar\omega$} \tag{variance of the pred. dist.}
\end{align*}

\subsection{GLMs}

\begin{enumerate}
	\item \emph{Linear regression} \textrightarrow{} regression ($\mathds{R}$)
	\item \emph{Logistic regression} \textrightarrow{} binary classification
	\item \emph{Poisson regression} \textrightarrow{} counter
	\item \emph{Multinomial regression} \textrightarrow{} multiclass classification
\end{enumerate}

\subsubsection{Multinomial regression}

We wish to create a GLM to predict one of $c$ classes: $\{ \omega_1, \omega_2, \cdots, \omega_c\}$

\begin{equation}
  x \in \begin{pmatrix}
    0 \\
    0 \\
    1 \\
    \vdots \\
    0
  \end{pmatrix} \Longleftrightarrow x \in \omega_3
  \tag {one-hot encoding}
\end{equation}

If we denote the probability of $x_l = 1$ by $\pi_l$

\begin{align*}
  P(x \mid \boldsymbol\pi) = \prod_{l=1}^c \pi_l^{x_l},\quad
  \boldsymbol{\pi} = (\pi_1, \pi_2, \cdots, \pi_c)^T,
  \pi_l \geq 0,\quad \sum_{l=1}^c \pi_l = 1
\end{align*}

\paragraph{Homework} Prove that:

\begin{align*}
  \sum_x P(x \mid \pi) &= 1 \\
  \mathds{E}[x \mid \boldsymbol\pi] &= \boldsymbol\pi
\end{align*}

\paragraph{Multinomial distribution}
The multinomial distribution is a generalization of the binomial distribution to more than two outcomes.

Given a number of observations of data $x_1, x_2, \cdots, x_n$. We define:

% and a set of probabilities $\pi_1, \pi_2, \cdots, \pi_k$ such that $\sum_{l=1}^k \pi_l = 1$, the probability of observing $x_1, x_2, \cdots, x_n$ is given by:

% \begin{equation}
  % P(x_1, x_2, \cdots, x_n \mid \boldsymbol\pi) = \frac{n!}{x_1! x_2! \cdots x_n!} \prod_{l=1}^k \pi_l^{x_l}
% \end{equation}

\begin{equation}
  m_l = \sum_{i=1}^n x_l^i \tag{number of observations of class $l$ in data}
\end{equation}

We consider now the joint distribution of the $m_1, m_2, \cdots, m_c$ conditional on the
parameters $\boldsymbol\pi$:

\begin{align*}
  \text{Mult}(\boldsymbol{m} \mid \boldsymbol\pi, n) &=
\underbrace{\binom{n}{m_1, m_2, \cdots, m_c}}_{\substack{\text{number of ways to distribute} \\ \text{$n$ observations among $m_1, m_2, \cdots, m_c$}}}
        \prod_{l=1}^c \pi_l^{m_l} \\
    \text{where } \sum_{l=1}^c m_l &= n \\
    \boldsymbol{m} &= (m_1, m_2, \cdots, m_c)^T
\end{align*}


\begin{enumerate}
  \item Obtaining the likelihood function $l = \mathcal{L}(\boldsymbol\pi)$.
  \item Maximize it \textrightarrow{} error function $- \ln \mathcal{L}(\boldsymbol\pi)$
  \item Softmax !
\end{enumerate}


%! TEX root = ../000-main.tex
% 2022-10-10

Now, we know that, using Bayes formula:

\begin{align*}
  P(\omega_i \mid x) &= \frac{P(x \mid \omega_i)P(\omega_i)}{\sum_{j=1}^c P(x \mid \omega_j)P(\omega_j)}, \; i=1 \dots c \\
                & \downarrow  a_i \coloneqq \ln \left\{ P(x \mid \omega_i)P(\omega_i) \right\} \\
                &= \frac{\exp(a_i)}{\sum_{j=1}^c exp(a_j)} \coloneqq f_i(x)
\end{align*}

This is called \emph{normalized exponential} or \emph{softmax}.

Suppose that $P(\omega_i \mid x) >> P(\omega_j  \mid x), \; i \neq j$
$\Rightarrow P(\omega_i \mid x) \leq 1, P(\omega_j \mid x) \geq 0, \; i \neq j$
$\sum_{j=1}^c P(x \mid \omega_j) = 1$.

Suppose further that we want to use a \emph{linear model}:

\begin{equation*}
  a_i = \omega^T \phi(x) = f_i(x), \; i = 1, \dots, c
\end{equation*}

% TODO: (homework)
% Show that the softmax reduces to the logistic regression when $c=2$.
%
% Hint: do the above with c=2 and simplify

Next we have to calculate the likelihood function (We'll use the log likelihood).

\begin{align*}
  \ell(\omega) \coloneqq \ln\mathcal{L}(\omega) &= \ln
  \prod_{i=1}^n \prod_{j=1}^c P\left(\omega_j \mid x^i\right)^{\tilde{y}_{ij}} \\
  &= \sum_{i=1}^n \sum_{j=1}^c \tilde{y}_{ij} \ln P(\omega_j \mid x^i) \\
  &= \sum_{i=1}^n \sum_{j=1}^c \tilde{y}_{ij} \ln \underbrace{f_j(x^i)}_{\hat{y}_{ij}} \\[1em]
  \tilde{y}_{ij} &\coloneqq \begin{cases}
    1 & \text{if } x_i \in \omega_j \\
    0 & \text{otherwise}
  \end{cases} \\
    E &\coloneqq -2\ln \mathcal{L} \tag{Error function / cross-entropy}
\end{align*}

We want to minimize the \emph{cross-entropy} $E$. To do so, we have to
calculate it numerically through Newton-Raphson \textrightarrow{} IRLS.

\begin{exercise}{}{} Consider a multiclassiflication problem for
$c$ classes defined by prior class probabilities
$\pi_j = P(\omega_j),\, j=1\dots, c$ and class densities
$p_j(x) = P(x\mid \omega_j)$. We are given training data
$D = \left\{ (x^i, y^i) \right\}_{i=1}^n$ where $y^i \in \{1, \dots, c\}$.

Obtain an estimation of the priors $\hat\pi_j$.


The solution is the proportion of data points labeled with class $\omega_j$.

\tcblower

\begin{align*}
  \hat\pi_j &= \frac{1}{n} \sum_{i=1}^n \mathds{1}\left\{ y^i = j \right\} \\
  &= \frac{1}{n} \sum_{i=1}^n \tilde{y}_{ij}
\end{align*}

\end{exercise}

\begin{exercise}{}{}
\begin{enumerate}
  \item Find estimators $\hat\pi_j f\sim\pi_j$
  \item How correct are they?
  \item Is there any other better estimator?
\end{enumerate}

\begin{hint}
  Look up statistical estimator statistic (bias, variance, consistency, efficiency)
\end{hint}
\end{exercise}

%! TEX root = ../000-main.tex
% 2022-10-24
\subsection{SVM exercises}

\begin{exercise}[breakable]{}{}
	Consider a hard-margin SVM. We know that for all Svs $x_i$:
	\begin{equation}
		\label{eq:svm:margin}
		\widehat\omega^T x^i + \widehat{b} = y^i
	\end{equation}

	\begin{enumerate}
		\item Obtain an expression for $\widehat{b}$ in therms of the
		      dual variable, $\widehat{\alpha}_i$.
		      \begin{align*}
			      \widehat\omega & = \sum_{i=1}^n \widehat\alpha_i y^i x^i                            \\
			      \widehat{b}    & = y^i - \left( \sum_{j=1}^n \widehat\alpha_j y^j x^j \right)^T x^i \\
			                     & = y^i - \sum_{j=1}^n \widehat\alpha_j y^j \left(x^j\right)^T x^i
		      \end{align*}

		      \begin{note}
			      In practice, we use more than one support vector:
			      \begin{equation*}
				      %\doublehat
				      \widehat b = \frac{1}{\left|S\right|} \sum_{i \in S}
				      \left(
				      y^i - \sum_{j=1}^n \widehat\alpha_j y^j \left(x^j\right)^T x^i
				      \right)
			      \end{equation*}
		      \end{note}

		\item Derive a simple expression for the margin of the dual.

      Using the expression for $\widehat b$ from the previous part, we can
      write the margin as:
		      \begin{align*}
            \widehat{b}   & = y^i - \sum_{j=1}^n \widehat\alpha_j y^j \left(x^j\right)^T x^i \label{eqstep:1}\tag{b} \\
            m &= \frac{2}{\lVert \widehat\omega \rVert} \tag{target}\\
            \sum_{i=1}^n \widehat\alpha_i y^i &= 0 \\
            \span\text{We take \cref{eqstep:1} and multiply by } (\widehat\alpha_iy^i): \\
            \widehat\alpha_iy^i \widehat b &= \widehat\alpha_iy^i y^i - \sum_{j=1}^n \widehat\alpha_iy^i \widehat\alpha_j y^j \left(x^j\right)^T x^i \\
            \cancelto{0}{\widehat\alpha_iy^i \widehat b} &=
            \widehat\alpha_i \cancelto{1}{y^i y^i} - \sum_{j=1}^n \widehat\alpha_iy^i \widehat\alpha_j y^j \left(x^j\right)^T x^i \\
            0 &= \sum_{i=1}^n \widehat\alpha_i
            - \sum_{i=1}^n \sum_{j=1}^n \widehat\alpha_i \widehat\alpha_j y^i y^j \left(x^j\right)^T x^i \\
            \sum_{i=1}^n \widehat\alpha_i &= \sum_{i=1}^n \sum_{j=1}^n \widehat\alpha_i \widehat\alpha_j y^i y^j \left(x^j\right)^T x^i = \left(\widehat\omega\right)^T \widehat\omega \\
            \sum_{i=1}^n \widehat\alpha_i &= \lVert \widehat\alpha \rVert_1 \\
            \widehat\alpha_i \geq 0 &\implies \lVert \widehat\alpha \rVert_1 = \lVert \widehat\omega \rVert^2 \\
            m &= \lVert \widehat\alpha \rVert_1^{-1/2}
		      \end{align*}
	\end{enumerate}
\end{exercise}

\begin{exercise}{}{}
	Consider again a hard-margin SVM in the dual form. Express
	the optimization problem as a minimization over a
	recognizable QP (quadratic problem), in vector-matrix form.
\end{exercise}

\begin{exercise}{}{}
  Consider this hypothesis space in $\mathds{R}$:
  \begin{equation*}
    \mathcal{F}_\text{SVM} \coloneqq \left\{
      x \to \text{sign}\left( \sin (\alpha x)\right),\; \alpha \in \mathds{R}
    \right\},\; x \in \mathds{R}
  \end{equation*}

  Calculate the VC-dimension of $\mathcal{F}_\text{SVM}$.

  \tcblower

  The VC-dimension is intrinsically related to the number
  of parameters.

  You choose $n \geq 2$

  \begin{align*}
    x^i \coloneqq 10^{-i},\; i \in \{1, \ldots, n\} \\
    \alpha_j \coloneqq \pi \left[
      1 + \sum_{i=1}^n \left( \frac{1}{2}(1 + ?)1-10^i \right)^i
    \right] \\
    y = (-1, +1, -1, +1, -1, +1, -1, +1, -1, -1,  +1, +1, -1, +1, +1)
  \end{align*}

\end{exercise}


\nocite{*}

\printindex\pagebreak
\printbibliography

\end{document}
